% Standalone source generated from article.tex and its local inputs.
% Regenerate with: python code/build_standalone.py
\documentclass[11pt,a4paper]{article}
\usepackage[margin=27mm,headheight=15pt]{geometry}
\usepackage[T1]{fontenc}
\usepackage{lmodern}
\usepackage{amsmath,amssymb,amsthm,mathtools,bm}
\usepackage{microtype,booktabs,longtable,array,enumitem,xcolor,xurl}
\usepackage{fancyhdr}
\definecolor{ink}{HTML}{18334A}
\definecolor{accent}{HTML}{176C73}
\usepackage[colorlinks=true,linkcolor=ink,citecolor=accent,urlcolor=accent]{hyperref}
\usepackage{bookmark}
\hypersetup{pdftitle={Mixed Spectral Identities: Harmonic Newton Series, Dougall Jets, Unequal Twists, and Stieltjes Collisions},pdfauthor={ProveIt research continuation, prepared with OpenAI}}
\numberwithin{equation}{section}
\newtheorem{theorem}{Theorem}[section]
\newtheorem{lemma}[theorem]{Lemma}
\newtheorem{proposition}[theorem]{Proposition}
\newtheorem{corollary}[theorem]{Corollary}
\theoremstyle{definition}
\newtheorem{definition}[theorem]{Definition}
\newtheorem{question}{Question}
\theoremstyle{remark}
\newtheorem{remark}[theorem]{Remark}
\newtheorem{example}[theorem]{Example}
\DeclareMathOperator{\Li}{Li}
\DeclareMathOperator{\FP}{FP}
\DeclareMathOperator{\CT}{CT}
\DeclareMathOperator{\Res}{Res}
\DeclareMathOperator{\sgn}{sgn}
\DeclareMathOperator{\supp}{supp}
\DeclareMathOperator{\PV}{PV}
\DeclareMathOperator{\Log}{Log}
\newcommand{\N}{\mathbb N}
\newcommand{\Z}{\mathbb Z}
\newcommand{\Q}{\mathbb Q}
\newcommand{\R}{\mathbb R}
\newcommand{\C}{\mathbb C}
\newcommand{\T}{\mathbb T}
\newcommand{\ii}{\mathrm i}
\newcommand{\dd}{\,\mathrm d}
\newcommand{\eps}{\varepsilon}
\newcommand{\fall}[2]{#1^{\underline{#2}}}
\newcommand{\stir}[2]{\genfrac{\{} {\}}{0pt}{}{#1}{#2}}
\newcommand{\coeff}[1]{\left[#1\right]}
\newcommand{\src}[1]{\nolinkurl{#1}}
\newcommand{\zetad}[2]{\zeta^{[#1]}(#2)}
\setlist{itemsep=3pt,topsep=5pt}
\setcounter{tocdepth}{2}
\emergencystretch=2em
\pagestyle{fancy}
\fancyhf{}
\fancyhead[L]{\small\color{ink}Mixed spectral identities}
\fancyhead[R]{\small ProveIt continuation}
\fancyfoot[C]{\thepage}
\renewcommand{\headrulewidth}{0.3pt}
\begin{document}
\hypersetup{pageanchor=false}
\begin{titlepage}
\centering
\vspace*{18mm}
{\large\color{accent}PROVEIT RESEARCH CONTINUATION\par}
\vspace{11mm}
{\LARGE\bfseries\color{ink}Mixed Spectral Identities\par}
\vspace{7mm}
{\Large Harmonic Newton Series, Dougall Jets,\\[4pt]
Unequal Twists, and Stieltjes Collisions\par}
\vspace{13mm}
{\large Detailed proofs and reproducible computations\par}
\vspace{5mm}
{10 October 2026\par}
\vspace{13mm}
\begin{minipage}{0.88\textwidth}
% BEGIN INLINED SOURCE: sections/abstract.tex
\textbf{Abstract.}
We continue four specific research directions in the ProveIt manuscript
and its six incoming reports at the pinned repository snapshot.
A normally convergent Newton expansion determines every spectral
derivative of the entire harmonic remainder, including at nonpositive
integers. Its shift jets are identified with classical multiple
Arakawa--Kaneko functions, and a two-shift boundary yields generalized
Stieltjes constants with finite repeated antiderivatives.

At the double residue zero of the centered Dougall family, we prove
termwise polynomial collapse, a divided generating identity, and an
all-order mixed-parameter table. Consequences include a weighted
trigamma square and harmonic squares equal to
$\zeta(2)\pm\zeta(3)$ and $3\zeta(4)-\zeta(5)$.
An additional spectral variable generates every Stieltjes order.
For unequal twists, we give a normally convergent Lerch expansion,
all-index finite-part convolutions, exact finite reductions, and
covariant antiderivatives.

For colliding Stieltjes products, we prove an all-factor
resonant-subset completion theorem and exhibit a finite anomaly:
the geometric triple collision with rates $(0,1,2)$ differs from
the merged finite part by $\tfrac12\zeta(2)\log2$.
We distinguish geometric collision limits, spectral restrictions,
and iterated Laurent constants. The article includes detailed
proofs, source corrections, reproducible diagnostics, and a focused
research agenda. The retained Gaussian $S_6$ and $S_8$ conjectures
are not resolved here.

% END INLINED SOURCE: sections/abstract.tex
\end{minipage}
\vfill
\begin{minipage}{0.88\textwidth}\small
Prepared with OpenAI for the ProveIt research programme.\\[4pt]
Repository snapshot:
\src{dcd95baeb7c0e914b138bddef6829c5b1d52b726}.\\[4pt]
Ordinary mathematical proofs are distinguished from exact finite
computations and numerical diagnostics. No proof-assistant
formalization or arithmetic-independence claim is made.
\end{minipage}
\end{titlepage}
\pagenumbering{roman}
\hypersetup{pageanchor=true}
\tableofcontents
\clearpage
\pagenumbering{arabic}
% BEGIN INLINED SOURCE: sections/01_scope.tex
\section{Scope, conventions, and source status}
\label{scope:sec}

\subsection{What this continuation establishes}

The task is to extend the exact-identity programme in the ProveIt
polylogarithm manuscript, not to reproduce results already supplied
by its incoming research reports. We inspected the canonical
manuscript and the six archives present in \src{docs/incoming}
at commit \src{dcd95baeb7c0e914b138bddef6829c5b1d52b726}
\cite{RepoMain,RepoEndpoint,RepoDougall,RepoTransport,RepoExact,
RepoCoincident,RepoTwisted}.
The mathematical work concentrates on four gaps explicitly left
by those sources. The following table records the resulting status.

\begin{longtable}{@{}>{\raggedright\arraybackslash}p{0.26\textwidth}>{\raggedright\arraybackslash}p{0.64\textwidth}@{}}
\toprule
Source target & Result and its precise scope\\
\midrule
\endhead
Higher spectral derivatives of the entire harmonic difference
\cite{RepoExact}
& Theorem~\ref{hn:thm:newton} and Corollary~\ref{hn:cor:jets}
give convergent formulas at every order. Theorem~\ref{hn:thm:AKbridge}
identifies all shift jets with classical multiple Arakawa--Kaneko
functions. A finite reduction of every general spectral derivative
to ordinary zeta derivatives is not asserted.\\[6pt]
Mixed Dougall jets at $b=c=1$, Question~3 of \cite{RepoTransport}
& Proposition~\ref{dg:polynomial} proves termwise collapse on the
integer residue grid. Theorems~\ref{dg:generator} and
\ref{dg:spectral} give every mixed parameter row and every
mixed spectral--Stieltjes derivative.\\[6pt]
Genuinely unequal twists, the first research question in \cite{RepoTwisted}
& Theorem~\ref{tw:tailtheorem} gives a convergent normal form at
arbitrary orders. The all-index formula \eqref{tw:alljets}
and finite scalar reductions retain the phase and endpoint
normalizations. The scalar classification is not an arithmetic
independence theorem for Fourier samples.\\[6pt]
Mixed triple collisions, Question~1 of \cite{RepoExact}
& Theorems~\ref{cl:collision} and \ref{cl:allfactor} give an
explicit triple collision expansion and the local completion for
any finite number of factors. Unqualified equality of the geometric
collision constant and the merged finite part is false. Theorem~\ref{cl:total-derivative} also resolves the one-derivative cancellation target.\\
\bottomrule
\end{longtable}

The results are proved below as analytic identities with specified
domains and normalizations. Exact finite calculations are supporting
certificates for selected algebraic steps. Numerical comparisons
test implementations and signs; they do not supply the infinite
identities or establish arithmetic independence. The supplied proofs
have not been formalized in a proof assistant.

\subsection{Classical antecedents and attribution}

The basic functions and summations are established mathematics.
We use the Hurwitz-zeta shift and derivative formulas, the
polygamma series, the Lerch transcendent, Euler's hypergeometric
transformation, and the Rogers--Dougall ${}_5F_4(1)$ summation
\cite{DLMFHurwitz,DLMFPolygamma,DLMFLerch,DLMFHypergeom,DLMFDougall}.
The multiple Arakawa--Kaneko functions were already defined in
\cite{AK1999}; related multiple and Hurwitz extensions are treated
in \cite{KT2018,CC2010}. The Hurwitz--Tornheim relation used in the
collision section has the antecedent \cite{EspinosaMoll}.
The centered resonance theorem and the entire harmonic subtraction
are results of the supplied reports. Where needed, their analytic
mechanisms are reproduced to make the new deductions reviewable
without reconstructing an entire earlier archive.

The contribution claimed here is the collection of deductions,
normal forms, convergence arguments, and corrections relative to
that inspected source set. We make no exhaustive priority claim
over the full special-functions literature. In particular, a
recognition of a classical Arakawa--Kaneko function is recorded as
an identification, not as the invention of a new special function.

\subsection{Notation and the meaning of finite part}

Write $(a)_n=\Gamma(a+n)/\Gamma(a)$ for a rising factorial,
$(a)^{\underline n}=a(a-1)\cdots(a-n+1)$ for a falling factorial,
and $[u^r]$ for coefficient extraction. The symbols
$\left\{\begin{smallmatrix}m\\j\end{smallmatrix}\right\}$ and
$\left[\begin{smallmatrix}m\\j\end{smallmatrix}\right]$ denote
Stirling numbers of the second kind and unsigned Stirling numbers
of the first kind, respectively. Empty products are $1$.
The ordinary generalized harmonic numbers are
\[
 H_n^{(r)}=\sum_{k=1}^n k^{-r},\qquad H_0^{(r)}=0,
 \qquad H_n=H_n^{(1)}.
\]
General shifts are defined explicitly in the relevant sections.
The letters $a,b,u$ are local parameters and are reintroduced
when their role changes.

We set $\psi=\Gamma'/\Gamma$, and $\psi^{(r)}$ denotes an
argument derivative. For zeta functions, $\zeta^{(k)}(s,x)$
denotes a derivative in the first, spectral variable; derivatives
in $x$ are written $\partial_x$. The generalized Stieltjes
constants use the convention
\begin{equation}\label{scope:stieltjes}
 \zeta(1+z,x)=\frac1z+
       \sum_{m=0}^{\infty}\frac{(-1)^m}{m!}\gamma_m(x)z^m,
 \qquad \gamma_m=\gamma_m(1),\quad \gamma=\gamma_0.
\end{equation}
In particular $\gamma_0(x)=-\psi(x)$.
Principal logarithms are used for powers in the right half-plane.
The phases for two-sided Fourier powers are fixed separately
in \eqref{tw:log}. A square at a complex shift is an algebraic
square, not an absolute-value square.

A series is \emph{normally convergent} on a parameter domain when
it converges absolutely and uniformly on each compact subset.
This is the hypothesis used for holomorphic differentiation of
infinite sums. The distribution-valued version means normal
convergence after pairing with each smooth test function, together
with the displayed uniform growth estimates on Fourier coefficients.

For a cutoff expansion in a positive ambient distance $\eta$,
$\FP$ retains the coefficient of $\eta^0(\log\eta)^0$ after
the divergent local terms are removed. At each singular grid
point in Section~\ref{cl:sec}, the cutoff is on the right in
the ambient coordinate $x$. This specifies a value; a change of
coordinate can change it. The symbol $\CT_\epsilon$ has the
analogous meaning for a geometric collision parameter $\epsilon$.
For multivariate meromorphic functions with poles on sums of
variables, an unqualified rectangular Laurent constant is not a
definition. A chamber, a spectral restriction, or the holomorphic
completion must first be specified.

\subsection{Reading the four proof mechanisms}

The harmonic argument first proves a compact-uniform estimate for
finite differences, then identifies their Mellin transform. This
order prevents the invalid differentiation of a formula known only
at discrete spectral values. The Dougall argument keeps the Gamma
ratios and the finite small-index terms intact until after the
subtractions have been differentiated. The unequal-twist argument
isolates finitely many modes before applying a binomial expansion,
so that the logarithm branches agree on the remaining tail.
The collision argument works with the complete local numerator of
each resonant subset, retaining finite terms as well as residues.
These are the analytic steps on which the extensions depend.

% END INLINED SOURCE: sections/01_scope.tex
\clearpage
% BEGIN INLINED SOURCE: sections/02_harmonic.tex
\section{A Newton expansion for the entire harmonic difference}
\label{hn:sec:newton}

The exact-identities report \cite{RepoExact} leaves open the explicit evaluation of the
higher spectral derivatives of its entire common-shift remainder.
Its finite Stirling formula has an infinite, normally convergent
extension. This extension gives every spectral derivative, continues
the shift farther as harmonic depth increases, and identifies all shift
jets at the unshifted base point with multiple Arakawa--Kaneko functions.
Those functions are classical; the contribution here is their precise
connection with the repository's entire difference and the convergent
spectral-derivative formula. No literature-wide priority is asserted.

\subsection{Definitions and the known entire remainder}

For $\Re a>0$, use principal powers and put
\[
 \begin{gathered}
 F_a(s,u)=\sum_{n\ge0}\frac{(a+u)_n}{(a)_n(n+a)^s},\qquad
 K_a(u)=\frac{\Gamma(a)\Gamma(1+u)}{\Gamma(a+u)},\\[3pt]
 D(s,u;a)=F_a(s,u)-K_a(u)F_1(s,u).
 \end{gathered}
\]
The first series initially converges for $\Re s>1+\Re u$.
The source proves that $D$ is entire in $s$, locally near $u=0$.
If
\[
 e_r(n;a)=[u^r]\prod_{\nu=0}^{n-1}\left(1+\frac{u}{a+\nu}\right),
 \qquad M_r(s;a)=\sum_{n\ge0}e_r(n;a)(n+a)^{-s},
\]
then its depth coefficients are
\[
 D_r(s;a)=[u^r]D(s,u;a)
 =M_r(s;a)-\sum_{j=0}^r[u^j]K_a(u)\,M_{r-j}(s;1).
\]
In particular $D_0(s;a)=\zeta(s,a)-\zeta(s)$, with its removable
singularity at $s=1$ understood. Each $D_r$ is entire in $s$.

We use the source's exact Mellin remainder
\begin{align}
 R_a(t,u)
 &=-\frac{a-1}{1+u}e^{-at}
       {}_2F_1(1,a+u;2+u;1-e^{-t}),\label{hn:eq:R}\\
 D(s,u;a)&=\frac1{\Gamma(s)}\int_0^\infty
                         t^{s-1}R_a(t,u)\,dt,\qquad \Re s>0.
 \label{hn:eq:mellin}
\end{align}
To check the normalization directly, the geometric generating function is
\[
 \sum_{n\ge0}\frac{(a+u)_n}{(a)_n}e^{-(n+a)t}
 =e^{-at}{}_2F_1(1,a+u;a;e^{-t})
 =K_a(u)\frac{e^{-t}}{(1-e^{-t})^{1+u}}+R_a(t,u).
\]
The Gauss connection formula about $z=1$ gives the last equality:
the regular coefficient is $-(a-1)/(1+u)$ and the singular
coefficient is $K_a(u)$; the remaining singular Gauss factor reduces
to $e^{(a-1)t}$ before multiplication by $e^{-at}$.
This initially holds for generic parameters, then throughout the
stated germ by continuation \cite{DLMFHypergeom}.
The kernel $R_a$ is analytic at zero and decays exponentially at infinity.
Taking Mellin transforms proves \eqref{hn:eq:mellin} first in the
original common half-plane and then for $\Re s>0$ by continuation. Taylor subtraction at
zero proves the entire continuation in $s$.

Define the finite exponential polynomial
\begin{equation}\label{hn:eq:Aj}
 A_j(s)=\sum_{\ell=1}^j(-1)^{j-\ell}\binom j\ell\ell^{1-s},
 \qquad j\ge1.
\end{equation}
It is entire in $s$. At a nonpositive integer,
\begin{equation}\label{hn:eq:Stirling}
 A_j(-m)=j!\left\{\begin{matrix}m+1\\j\end{matrix}\right\},
 \qquad m\ge0,
\end{equation}
and therefore it vanishes for $j>m+1$.

\subsection{Normal convergence, including at nonpositive integers}

\begin{lemma}[Finite-difference Mellin estimate]\label{hn:lem:Aj}
For $\Re s>1-j$,
\begin{equation}\label{hn:eq:Aintegral}
 A_j(s)=\frac{(-1)^{j-1}j}{\Gamma(s)}
       \int_0^\infty t^{s-1}e^{-t}(1-e^{-t})^{j-1}\,dt.
\end{equation}
At zeros of $1/\Gamma(s)$ this is its analytic product with the
convergent integral. For every compact $S\subset\mathbb C$ there
are constants $C,P$ such that
\begin{equation}\label{hn:eq:Abound}
 \sup_{s\in S}|A_j(s)|\le C(1+\log(j+1))^P\qquad(j\ge1).
\end{equation}
Every fixed spectral derivative obeys a bound of the same form,
with constants depending on its order and $S$.
\end{lemma}

\begin{proof}
For $\Re s>0$, expand the finite binomial power and use
$j\binom{j-1}{\ell-1}=\ell\binom j\ell$.
Termwise Laplace integration gives \eqref{hn:eq:Aintegral}.
The integrand at zero is $O(t^{\Re s+j-2})$, so its integral is
holomorphic on $\Re s>1-j$. Analytic continuation gives the formula
on that larger half-plane. Thus it holds near any prescribed
nonpositive integer as soon as $j$ is sufficiently large.

Write $\sigma_-=\min_{s\in S}\Re s$ and
$\sigma_+=\max_{s\in S}\Re s$. Choose an integer $N$ such that
$\sigma_-+N-1>0$, and consider $j\ge N$. Put $q=1-e^{-1}<1$.
On $0<t\le1$,
\[
 (1-e^{-t})^{j-1}\le t^{N-1}q^{j-N},
\]
so the integral on this interval, including its factor $j$, is at
most $j q^{j-N}/(\sigma_-+N-1)$. This is exponentially small.
On $t\ge1$, put $P=\max(0,\sigma_+-1)$. The function
\[
                 p_j(t)=je^{-t}(1-e^{-t})^{j-1}
\]
is a probability density on $(0,\infty)$. With
$L_j=1+2\log(j+1)$, the integral on $[1,L_j]$ is bounded by $L_j^P$.
On $[L_j,\infty)$ use $p_j(t)\le je^{-t}$ and
\[
 j\int_{L_j}^\infty t^Pe^{-t}\,dt
 \le C_Pje^{-L_j}(1+L_j)^P
 \le C'_P(1+\log(j+1))^P.
\]
The reciprocal gamma function is bounded on $S$. The finitely many
remaining entire functions $A_j$ are harmless. Cauchy's estimate
on a slightly enlarged compact proves the derivative assertion.
\end{proof}

\begin{lemma}[Uniform binomial estimate]\label{hn:lem:binomial}
For a compact set $A\subset\mathbb C$,
\[
 \left|\binom{a-1}{j}\right|\le C_A j^{-\Re a}
 \qquad(a\in A,\ j\ge1).
\]
\end{lemma}
\begin{proof}
Use $(-1)^j\binom{a-1}{j}=\prod_{\nu=1}^j(1-a/\nu)$.
For $\nu>2\max_{a\in A}|a|$,
$\log|1-a/\nu|=-\Re a/\nu+O_A(\nu^{-2})$, uniformly.
Summing and bounding the initial finite factors proves the claim,
including when $a$ passes through a positive integer.
\end{proof}

\begin{theorem}[Normally convergent Newton expansion]\label{hn:thm:newton}
For $\Re a>0$ and $u\notin\{-1,-2,\ldots\}$,
\begin{equation}\label{hn:eq:newton}
 \boxed{\displaystyle
 D(s,u;a)=-\sum_{j=1}^\infty
                \binom{a-1}{j}\frac{A_j(s)}{u+j}.}
\end{equation}
The series converges normally on compact subsets of the indicated
domain with $s\in\mathbb C$. It extends the depth germ to a
meromorphic function of $u$, with at most simple poles at negative
integers, and
\begin{equation}\label{hn:eq:uresidue}
 \mathop{\rm Res}_{u=-j}D(s,u;a)=-\binom{a-1}{j}A_j(s).
\end{equation}
These poles are removable when the indicated residues vanish.

For every $r\ge0$,
\begin{equation}\label{hn:eq:depthnewton}
 \boxed{\displaystyle
 D_r(s;a)=(-1)^{r+1}\sum_{j=1}^\infty
          \binom{a-1}{j}\frac{A_j(s)}{j^{r+1}}.}
\end{equation}
This depth series converges normally on
\begin{equation}\label{hn:eq:shift-domain}
                  s\in\mathbb C,\qquad \Re a>-r,
\end{equation}
and supplies a jointly holomorphic continuation there.
\end{theorem}

\begin{proof}
On a compact $u$-set avoiding the negative integers,
$|u+j|^{-1}\le C/j$. The lemmas bound the summands by
$Cj^{-1-\delta}(1+\log(j+1))^P$, where
$\delta=\min\Re a>0$. In the depth series the same proof uses
$\delta=\min(\Re a+r)>0$. This proves normal convergence and
permits arbitrary fixed derivatives in $s$ and $a$.

To identify the sum, put $w=1-e^{-t}$. Euler's hypergeometric
transformation \cite{DLMFHypergeom} and then its series give
\begin{align*}
 R_a(t,u)
 &=-\frac{a-1}{1+u}e^{-t}
       {}_2F_1(2-a,1+u;2+u;w)\\
 &=e^{-t}\sum_{j=1}^\infty
        (-1)^j\binom{a-1}{j}\frac{j}{u+j}w^{j-1}.
\end{align*}
Initially $u$ is near zero and $0<w<1$. For $\Re s>0$ this
series can be integrated absolutely against $t^{s-1}\,dt$:
apply \eqref{hn:eq:Aintegral} with $s$ replaced by $\Re s$,
and the estimates already proved. Its Mellin transform is
\[
 \frac1{\Gamma(s)}\int_0^\infty t^{s-1}R_a(t,u)\,dt
 =-\sum_{j\ge1}\binom{a-1}{j}\frac{A_j(s)}{u+j}.
\]
It equals $D$ in a common domain; both sides are entire in $s$.
The normal convergence gives the continuation and residues in $u$.
Expand $(u+j)^{-1}$ on $|u|<1$ and extract depth coefficients.
The enlarged shift domain follows from direct normal convergence,
not interpolation from integer shifts.
\end{proof}

\subsection{All spectral derivatives}

\begin{corollary}[Logarithmic finite-difference formula]\label{hn:cor:jets}
For $m,k,r\ge0$,
\begin{align}
 \partial_s^kD(-m,u;a)
 &=(-1)^{k+1}\sum_{j\ge1}\frac{\binom{a-1}{j}}{u+j}
       \sum_{\ell=1}^j(-1)^{j-\ell}\binom j\ell
                     \ell^{m+1}(\log\ell)^k,
 \label{hn:eq:spectral-full}\\
 \partial_s^kD_r(-m;a)
 &=(-1)^{r+k+1}\sum_{j\ge1}\frac{\binom{a-1}{j}}{j^{r+1}}
       \sum_{\ell=1}^j(-1)^{j-\ell}\binom j\ell
                     \ell^{m+1}(\log\ell)^k.
 \label{hn:eq:spectral-depth}
\end{align}
The first formula is normally convergent for $\Re a>0$, $u$ away
from the negative integers; the second for $\Re a>-r$.
At $k=0$ use $(\log1)^0=1$.
\end{corollary}
\begin{proof}
Differentiate the normally convergent series using the compact
spectral-derivative estimates.
\end{proof}

At $k=0$ this is exactly the source's finite Stirling formula,
\[
 D(-m,u;a)=-\sum_{j=1}^{m+1}
       \left\{\begin{matrix}m+1\\j\end{matrix}\right\}
       \frac{(a-1)^{\underline j}}{u+j}.
\]
For $k\ge1$ the finite difference generally no longer vanishes at
large $j$. The answer is an absolutely convergent infinite identity;
no finite reduction to polygamma functions at arbitrary depth is
claimed. A finite value formula valid only at discrete values
of $s$ cannot itself be differentiated in $s$.

For clarity, if $[\epsilon^k]$ denotes a Laurent coefficient at
$s=-m$, the identity defining $D_r$ gives, for $k\ge0$,
\begin{align}
 [\epsilon^k]M_r(-m+\epsilon;a)
 ={}&\sum_{j=0}^r[u^j]K_a(u)\,
                 [\epsilon^k]M_{r-j}(-m+\epsilon;1)\notag\\
 &+\frac1{k!}\partial_s^kD_r(-m;a).
 \label{hn:eq:regular-reconstruction}
\end{align}
Thus the theorem determines every coefficient of the entire
shift correction. It does not by itself reduce the unshifted
height-one Laurent coefficients, or the entire outer-shift family
of the earlier report, to a finite basis of ordinary zeta derivatives.


Depth-zero special cases give useful independent checks:
\begin{align}
 \log\Gamma(a)
 &=\sum_{j\ge1}\frac{\binom{a-1}{j}}j
      \sum_{\ell=1}^j(-1)^{j-\ell}\binom j\ell\ell\log\ell,
 \label{hn:eq:loggamma}\\
 \gamma_k(a)-\gamma_k(1)
 &=-\sum_{j\ge1}\frac{\binom{a-1}{j}}j
      \sum_{\ell=1}^j(-1)^{j-\ell}\binom j\ell(\log\ell)^k.
 \label{hn:eq:Stieltjes-Newton}
\end{align}
The Stieltjes convention is
$\zeta(1+z,a)=z^{-1}+\sum_{k\ge0}(-1)^k\gamma_k(a)z^k/k!$.
The first identity follows from Lerch's formula \cite{DLMFHurwitz} and the second
from the regular expansion of $D_0$ at $s=1$. These are
specializations of a Newton/Hasse type calculus, with no claim
that the classical log-gamma or Stieltjes expansions are new.

There is a finite reduction at every nonnegative integer
\emph{depth-generating parameter} $u=q$. For $q\ge1$ and initially $\Re a>0$,
\begin{equation}\label{hn:eq:integeru}
 D(s,q;a)=\frac1{(a)_q}\sum_{\ell=1}^q
      \left[\begin{matrix}q\\\ell\end{matrix}\right]
            \{\zeta(s-\ell,a)-\zeta(s-\ell)\},
\end{equation}
where brackets denote unsigned Stirling numbers of the first kind.
Indeed $(a+q)_n/(a)_n=(n+a)_q/(a)_q$, and expand the rising
factorial polynomial for both $F_a$ and $K_aF_1$.
All spectral derivatives follow by replacing each zeta by its
order derivative. In particular
\[
 \partial_s^kD(-m,1;a)
   =\frac{\zeta^{(k)}(-m-1,a)-\zeta^{(k)}(-m-1)}a.
\]
Every zeta difference in \eqref{hn:eq:integeru} is entire,
including at its removable principal pole.

\subsection{The Arakawa--Kaneko boundary and every shift jet}

For a positive integer multi-index
$\boldsymbol{k}=(k_1,\ldots,k_p)$ use the descending-index convention
\[
 \operatorname{Li}_{\boldsymbol{k}}(z)
 =\sum_{n_1>\cdots>n_p\ge1}
       \frac{z^{n_1}}{n_1^{k_1}\cdots n_p^{k_p}}.
\]
Define the classical multiple Arakawa--Kaneko function by
\begin{equation}\label{hn:eq:AK}
 \Xi(\boldsymbol{k};s)=\frac1{\Gamma(s)}\int_0^\infty
       \frac{t^{s-1}}{e^t-1}
              \operatorname{Li}_{\boldsymbol{k}}(1-e^{-t})\,dt,
 \qquad \Re s>0.
\end{equation}
Write $\Xi_r(s)=\Xi((r);s)$ for $r\ge1$.
Arakawa--Kaneko \cite{AK1999} already define the multiple function.
They use ascending summation indices, so
$\Xi(r+1,\{1\}^{p-1};s)$ here corresponds to their
$\xi(1,\ldots,1,r+1;s)$; the vector is reversed.
At depth one there is no reversal ambiguity, and
\[
                         \Xi_1(s)=s\zeta(s+1),
\]
with a removable singularity at $s=0$.

For completeness all functions in \eqref{hn:eq:AK} are entire in $s$.
The polylogarithm is analytic at $z=0$ and begins at degree $p$,
so its kernel is analytic at $t=0$, beginning at order $p-1$.
At infinity its growth before multiplication by $(e^t-1)^{-1}$
is at most polynomial: since all indices are positive, its
positive terms are bounded by those of
$\operatorname{Li}_{1,\ldots,1}(1-e^{-t})=t^p/p!$.
Taylor subtraction at zero, followed by the zeros of $1/\Gamma(s)$,
therefore proves entireness. This is also the continuation in
the classical theory \cite{AK1999,KT2018}.

\begin{theorem}[Boundary values and mixed shift jets]\label{hn:thm:AKbridge}
For $r\ge1$, the continuation in \eqref{hn:eq:shift-domain} gives
\begin{equation}\label{hn:eq:a0}
                      D_r(s;0)=(-1)^r\Xi_r(s).
\end{equation}
For every $r\ge0$ and $p\ge1$,
\begin{equation}\label{hn:eq:a1-jets}
 \boxed{\displaystyle
 \left.\partial_a^pD_r(s;a)\right|_{a=1}
  =(-1)^{r+p}p!\,
          \Xi(r+1,\underbrace{1,\ldots,1}_{p-1};s).}
\end{equation}
The identities hold for every $s\in\mathbb C$, and arbitrary
fixed spectral derivatives can be applied to both sides.
Consequently
\begin{equation}\label{hn:eq:AK-Taylor}
 D_r(s;a)=(-1)^r\sum_{p=1}^\infty
       (1-a)^p\Xi(r+1,\{1\}^{p-1};s),
 \qquad |a-1|<r+1.
\end{equation}
The series is locally normal jointly with $s$.
\end{theorem}

\begin{proof}
For $a=0$, $\binom{-1}{j}=(-1)^j$. Substitute into the normally
convergent depth series and apply \eqref{hn:eq:Aintegral}:
\[
 D_r(s;0)=\frac{(-1)^r}{\Gamma(s)}\int_0^\infty
     t^{s-1}e^{-t}\sum_{j\ge1}\frac{(1-e^{-t})^{j-1}}{j^r}\,dt.
\]
For $\Re s>0$, termwise integration is justified absolutely.
The sum is $\operatorname{Li}_r(w)/w$, where $w=1-e^{-t}$,
giving \eqref{hn:eq:a0}. Continue in $s$.

For the shift jets let $e_j(N)$ be the elementary symmetric
polynomial of degree $j$ in $1,1/2,\ldots,1/N$, with $e_0(N)=1$
and $e_j(N)=0$ for $j>N$. The finite product
\[
 \binom{a-1}{n}=\frac{(-1)^{n-1}(a-1)}n
          \prod_{\nu=1}^{n-1}\left(1-\frac{a-1}{\nu}\right)
\]
gives
\[
 \left.\partial_a^p\binom{a-1}{n}\right|_{a=1}
      =\frac{(-1)^{n-p}p!}{n}e_{p-1}(n-1).
\]
Normal convergence permits insertion into
\eqref{hn:eq:depthnewton}. The Mellin integral then contains
\[
 \sum_{n\ge1}\frac{w^{n-1}}{n^{r+1}}e_{p-1}(n-1)
 =\frac1w\operatorname{Li}_{r+1,\{1\}^{p-1}}(w),
\]
with overall sign $(-1)^{r+p}p!$. Absolute integration follows
from positive coefficients and the logarithmic growth
$e_{p-1}(n-1)=O((1+\log n)^{p-1})$.
This proves \eqref{hn:eq:a1-jets} on $\Re s>0$, then everywhere.

Finally $D_r(s;1)=0$, and $D_r$ is holomorphic for $\Re a>-r$.
Its Taylor series centered at $1$ therefore converges throughout
$|a-1|<r+1$. Cauchy estimates on compact $s$-sets prove joint
local normality.
\end{proof}

The depth-zero case of \eqref{hn:eq:a1-jets} reduces to
\[
 \Xi(\underbrace{1,\ldots,1}_{p};s)
        =\frac{(s)_p}{p!}\zeta(s+p).
\]
The boundary formula \eqref{hn:eq:a0}, however, excludes $r=0$.
The original domain does not include $a=0$, and
\[
 D_0(s;a)=a^{-s}+\zeta(s,a+1)-\zeta(s)
\]
has no jointly holomorphic extension there for general $s$.
A fictitious $\Xi_0$ with a nondecaying Mellin kernel cannot
remove this obstruction.

Some exact examples, with $\gamma_k=\gamma_k(1)$, are
\begin{align*}
 D_1(s;0)&=-s\zeta(s+1),\\
 D_1(0;0)&=-1,&
 \partial_sD_1(0;0)&=-\gamma,&
 \partial_s^2D_1(0;0)&=2\gamma_1,\\
 \partial_s^kD_1(0;0)&=(-1)^k k\gamma_{k-1}\quad(k\ge1),\\
 \partial_sD_1(-1;0)&=\frac12-\frac12\log(2\pi),\\
 \partial_sD_1(-2;0)&=\frac1{12}+2\zeta'(-1),\\
 \partial_aD_1(s;1)&=\Xi_2(s),&
 \partial_a^2D_1(s;1)&=-2\Xi(2,1;s).
\end{align*}
They are identities for the continued remainder, rather than
ordinary sums of its divergent defining series.

More generally, for an integer $N\ge0$ and depth $r\ge N+1$,
\begin{equation}\label{hn:eq:negative-shifts}
 D_r(s;-N)=\frac{(-1)^r}{N!}\sum_{h=0}^N
     \left[\begin{matrix}N+1\\h+1\end{matrix}\right]\Xi_{r-h}(s).
\end{equation}
Indeed $\binom{-N-1}{j}=(-1)^j\binom{N+j}{N}$ and
\[
 \binom{N+j}{N}=\frac1{N!}\prod_{\nu=1}^N(j+\nu)
   =\frac1{N!}\sum_{h=0}^N
       \left[\begin{matrix}N+1\\h+1\end{matrix}\right]j^h.
\]
Insert this polynomial into the convergent Mellin series. The
condition $r>N$ makes all polylogarithm indices positive and
places the shift inside $\Re a>-r$. For example
\[
                   D_2(s;-1)=\Xi_2(s)+s\zeta(s+1).
\]

Specializing the convergent Taylor expansion to $a=0$ also gives
$\Xi_r(s)=\sum_{p\ge1}\Xi(r+1,\{1\}^{p-1};s)$ for $r\ge1$,
normally on compact $s$-sets. Its elementary polylogarithmic
explanation is $\sum_{p\ge1}e_{p-1}(n-1)=n$.
This is recorded as a consequence of the shift identification,
with no new-priority assertion for the height-one summation itself.

\subsection{A boundary corollary for generalized Stieltjes constants}

The two-shift bridge in the source admits the same proof. Write
$v=b-a$, $c=1+v$, and initially assume $\Re a,\Re b,\Re c>0$.
Define
\[
 D_{a,b}(s,u)=\sum_{n\ge0}\frac{(a+u)_n}{(a)_n(n+b)^s}
   -K_a(u)\sum_{n\ge0}\frac{(1+u)_n}{n!(n+c)^s},
 \qquad D_{a,b,r}(s)=[u^r]D_{a,b}(s,u).
\]
Its entire Mellin remainder is $e^{-vt}R_a(t,u)$.
Put
\[
 A_j(s;v)=\sum_{\ell=1}^j(-1)^{j-\ell}\binom j\ell
                     \ell(\ell+v)^{-s},\qquad \Re v>-1.
\]
The Laplace calculation gives
\[
 A_j(s;v)=\frac{(-1)^{j-1}j}{\Gamma(s)}
       \int_0^\infty t^{s-1}e^{-(1+v)t}(1-e^{-t})^{j-1}\,dt,
 \qquad \Re s>1-j.
\]
Then
\begin{equation}\label{hn:eq:twodepth}
 D_{a,b,r}(s)=(-1)^{r+1}\sum_{j\ge1}\binom{a-1}{j}
                           \frac{A_j(s;b-a)}{j^{r+1}},
\end{equation}
normally on the connected domain
\begin{equation}\label{hn:eq:two-domain}
              \Re(1+b-a)>0,\qquad \Re b>-r,\qquad s\in\mathbb C.
\end{equation}
This supplies a continuation without a restriction on $\Re a$.

Here is the convergence detail needed for this enlarged domain.
On compact $s,v$ sets with $\Re v>-1$, and $0<\epsilon<\min\Re(1+v)$ sufficiently small, the modified Mellin estimate is
$|A_j(s;v)|\le C_\epsilon j^{-\Re v+\epsilon}$.
On $(0,1)$ the previous proof is unchanged; on $t\ge1$ absorb
the necessary power of $t$ into $C_\epsilon e^{\epsilon t}$ and use
\[
 j\int_0^\infty e^{-(1+\Re v-\epsilon)t}
                  (1-e^{-t})^{j-1}\,dt
 =jB(1+\Re v-\epsilon,j)
 \le C_\epsilon j^{-\Re v+\epsilon}.
\]
The last bound follows directly by bounding
$(1-x)^{j-1}$ by $e^{-(j-1)x}$ in the beta integral.
Since $a+v=b$, choose $\epsilon$ below half of
$\min(\Re b+r)$ on the compact under consideration. The series
is then bounded by a summable power
$Cj^{-1-(\Re b+r)+\epsilon}$. Identification uses the same
kernel expansion as before. All fixed spectral derivatives are
allowed by normal convergence.

For $r\ge1$, $\Re x>0$, define the classical Hurwitz extension
\[
 \Xi_r(s,x)=\frac1{\Gamma(s)}\int_0^\infty
       \frac{e^{-xt}}{1-e^{-t}}t^{s-1}
                        \operatorname{Li}_r(1-e^{-t})\,dt.
\]
This convention agrees with Coppo--Candelpergher \cite{CC2010}, and
$\Xi_r(s,1)=\Xi_r(s)$. It is entire in $s$ by the same Taylor
subtraction argument. The point $a=0$, $b=x-1$ lies in
\eqref{hn:eq:two-domain} for $r\ge1$, so
\begin{equation}\label{hn:eq:generalized-boundary}
                  D_{0,x-1,r}(s)=(-1)^r\Xi_r(s,x).
\end{equation}
At depth one, $\operatorname{Li}_1(1-e^{-t})=t$ gives
\begin{equation}\label{hn:eq:Hurwitz-boundary}
 D_{0,x-1,1}(s)=-s\zeta(s+1,x)=\partial_x\zeta(s,x).
\end{equation}
The right side is entire in $s$, including at zero.

\begin{corollary}[Generalized Stieltjes boundary jets]\label{hn:cor:Stieltjes}
For $\Re x>0$ and $k\ge1$,
\begin{equation}\label{hn:eq:generalized-Stieltjes}
 \boxed{\displaystyle
 \left.\partial_s^kD_{0,x-1,1}(s)\right|_{s=0}
      =(-1)^k k\,\gamma_{k-1}(x).}
\end{equation}
The zeroth derivative is $-1$. For $m\ge1$,
\[
 \left.\partial_s^kD_{0,x-1,1}(s)\right|_{s=-m}
 =m\zeta^{(k)}(1-m,x)-k\zeta^{(k-1)}(1-m,x),
\]
with the last term omitted when $k=0$.
At $k=1,m=0$ the boundary function is $\psi(x)$, and its
antiderivative is $\log\Gamma(x)$.
\end{corollary}
\begin{proof}
Insert the defining Laurent expansion of the Hurwitz zeta
function in \eqref{hn:eq:Hurwitz-boundary}. The factor $s$
removes its pole before coefficient extraction.
\end{proof}

Every one of these boundary jets also has a finite repeated
antiderivative. For $m,k\ge0$ and $q\ge1$, define coefficients by
\[
 \prod_{\nu=1}^{q-1}\left(1-\frac z{m+\nu}\right)^{-1}
     =\sum_{h\ge0}h_h^{[m,q-1]}z^h
\]
(the product is $1$ at $q=1$). Then
\begin{equation}\label{hn:eq:boundary-primitives}
 \mathcal P_{m,k}^{(q)}(x)=\frac{k!}{(m+1)_{q-1}}
       \sum_{j=0}^k\frac{h_{k-j}^{[m,q-1]}}{j!}
                    \zeta^{(j)}(1-m-q,x)
\end{equation}
satisfies
\[
 \frac{d^q}{dx^q}\mathcal P_{m,k}^{(q)}(x)
       =\left.\partial_s^kD_{0,x-1,1}(s)\right|_{s=-m}.
\]
Indeed a $q$-fold primitive of $-s\zeta(s+1,x)$ is
$\zeta(s-q+1,x)/(1-s)_{q-1}$. Its denominator is nonzero near
$s=-m$, and differentiating $k$ times in $s$ proves the formula.
The coefficients are complete homogeneous symmetric polynomials
in $1/(m+1),\ldots,1/(m+q-1)$, equivalently Bell polynomials in
generalized harmonic-number differences. Other $q$-fold
primitives differ by a polynomial of degree less than $q$.

\subsection{Scope and further questions}

No substantive error was found in the inspected fully-shifted
and outer-shift harmonic arguments. Their restriction of the
finite coefficient rule to the principal part and finite part
is correct. The new formulas identify the missing analytic
information in higher spectral derivatives. Three extension
requirements are explicit: a discrete finite value formula
cannot be differentiated in $s$; depth and spectral manipulations
must use the already entire remainder instead of disregarding
the moving polar divisors of $F_a$; and the common-shift $a=0$
Arakawa--Kaneko formula requires positive depth.

Further questions include identifying finite combinations of
spectral jets that reduce to ordinary Hurwitz-zeta derivatives;
determining the exact shift singularities on $\Re a=-r$;
developing stable evaluation of the finite differences, since
naive binomial summation loses precision; comparing the present
shift deformation systematically with existing generalized
Arakawa--Kaneko families; and seeking multiple-zeta or
iterated-integral identities for mixed spectral and shift
derivatives beyond the standard integer specializations.


% END INLINED SOURCE: sections/02_harmonic.tex
\clearpage
% BEGIN INLINED SOURCE: sections/03_dougall.tex
\section{Mixed Dougall derivatives at a double residue zero}
\label{dg:section}

The centered resonance theorem already evaluates the minimal
subtraction at every nonpositive integer excess parameter. Its
residue-zero loci contain further identities: after restriction to
some loci the remainder vanishes term by term, while derivatives
transverse to those loci produce nonzero sums. We make this statement
precise, give a finite formula at every mixed order, and evaluate a
weighted square of trigamma differences for an arbitrary positive
shift. An additional spectral variable gives an all-order generating
function for generalized Stieltjes constants.

The underlying summation is the classical Rogers--Dougall
${}_5F_4(1)$ formula, DLMF 16.4.9~\cite{DLMFDougall}. The summation itself and the
centered subtraction theorem are antecedents, not novelty claims~\cite{RepoDougall,RepoTransport}.
The formulas below resolve the explicit mixed-parameter target
at $b=c=1$ in the incoming coordinate-transport report. No claim of
global literature priority is required for these deductions.

\subsection{Exact polynomial collapse on the integer grid}

Use the centered notation
\[
 \eta=a/2,\quad x_n=n+\eta,\quad d_0=1+a-b-c,
\]
and
\[
 r_n(u)=x_n\frac{\Gamma(n+a)\Gamma(n+b)\Gamma(n+c)
                     \Gamma(n+d_0-u)}
 {\Gamma(n+1)\Gamma(n+1+a-b)\Gamma(n+1+a-c)
                     \Gamma(n+b+c+u)}.
\]
Reciprocal Gamma factors are interpreted as entire functions, including
at their zeros. Put
\[
 (\alpha_1,\alpha_2,\alpha_3,\alpha_4)
       =(a/2,b-a/2,c-a/2,d_0-u-a/2),
\]
and define finite polynomial data by
\begin{align}
 h_j(u)&=-\frac{\sum_{i=1}^4B_{2j+1}(\alpha_i)}{j(2j+1)},\qquad j\geq1,
       \notag\\
 C_0(u)&=1,\qquad
 jC_j(u)=\sum_{k=1}^jkh_k(u)C_{j-k}(u).
 \label{dg:centered-recurrence}
\end{align}
Here $B_j(z)$ denotes the Bernoulli polynomial with $B_1(z)=z-1/2$.
These are exactly the centered asymptotic coefficients:
\[
 r_n(u)\sim x_n^{-1-2u}\sum_{j\geq0}C_j(u)x_n^{-2j},\qquad C_0=1.
\]
At $u=-N$, $N\geq1$, define the normally convergent bracket
\[
 F_{N,n}(a,b,c)=r_n(-N)-\sum_{j=0}^N C_j(-N)x_n^{2N-1-2j}.
\]

\begin{lemma}[Centered summation and residue transport: antecedent]
\label{dg:antecedent}
Suppose $a,b,c>0$ and $d_0>0$. For $K\geq1$ and
$-K<\Re u<d_0$, the normally convergent series
\begin{align}
 R_K(u)&=\sum_{n\geq0}\left\{r_n(u)-
       \sum_{j=0}^{K-1}C_j(u)x_n^{-1-2u-2j}\right\}\notag\\
 &=G(u)-\sum_{j=0}^{K-1}C_j(u)\zeta(1+2u+2j,a/2),\notag\\
 G(u)&=\frac{\Gamma(b)\Gamma(c)\Gamma(d_0-u)\Gamma(u)}
          {2\Gamma(d_0)\Gamma(b+u)\Gamma(c+u)}.
 \label{dg:antecedent-master}
\end{align}
The combined right side is holomorphic throughout this strip.
For $N\geq0$, let
\[
 \rho_N=\frac{(-1)^N}{N!}(d_0)_N(1-b)_N(1-c)_N.
\]
Then
\begin{align}
 C_N(-N)&=\rho_N,\notag\\
 R_{N+1}(-N)&=\rho_N\left\{\psi(a/2)+
 \frac{\psi(N+1)-\psi(b)-\psi(c)-\psi(d_0+N)}2\right\}.
 \label{dg:antecedent-closed}
\end{align}
All fixed spectral and parameter derivatives may be taken in the
single bracket in \eqref{dg:antecedent-master}.
\end{lemma}
\begin{proof}
At $x_n=n+a/2$, the denominator shifts are $1-\alpha_i$, so
\[
 r_n(u)=x_n\prod_{i=1}^4
                   \frac{\Gamma(x_n+\alpha_i)}{\Gamma(x_n+1-\alpha_i)}.
\]
The finite Stirling--Bernoulli expansion~\cite{DLMFGammaAsymptotic} and Bernoulli reflection
give, uniformly on compact parameter sets,
\[
 r_n(u)=x_n^{-1-2u}\left\{\sum_{j=0}^{K-1}C_j(u)x_n^{-2j}
                                     +O(n^{-2K})\right\}.
\]
Only even inverse powers occur; exponentiating the finite logarithmic
expansion gives \eqref{dg:centered-recurrence}. Thus the bracket is
summable throughout $-K<\Re u<d_0$. Cauchy's formula on a slightly
larger compact set proves the same for each derivative, with fixed
logarithmic powers allowed. In $0<\Re u<d_0$, the classical
Rogers--Dougall formula~\cite{DLMFDougall} sums $r_n(u)$ to $G(u)$:
the upper parameters are $a,1+a/2,b,c,d_0-u$, the lower parameters
are $a/2,1+a-b,1+a-c,b+c+u$, and the conventional excess is $2u$.
Separate summation of the counterterms gives
\eqref{dg:antecedent-master} there. Holomorphic continuation proves
the full-strip identity and removability of its apparent poles.

At $u=-N$ its Gamma residue is $\rho_N/2$. The only polar Hurwitz
counterterm has index $N$ and residue $C_N(-N)/2$. Equality of
residues gives $C_N(-N)=\rho_N$; since these are polynomial
coefficients, this equality extends through every residue zero.

For the constant term use
\[
 V=2\partial_a+\partial_b+\partial_c-\partial_u,
 \qquad W=\partial_a+\tfrac12\partial_b+\tfrac12\partial_c.
\]
One has $V\alpha_1=1$ and $V\alpha_i=0$ for $i=2,3,4$.
Consequently
\[
 Vh_k(u)=-\frac{B_{2k}(a/2)}k=2\zeta(1-2k,a/2),
 \qquad VC_j(u)=2\sum_{k=1}^j\zeta(1-2k,a/2)C_{j-k}(u).
\]
The second formula follows by differentiating the formal exponential
$\sum_jC_jy^j=\exp(\sum_kh_ky^k)$; every extracted coefficient is
finite. At $u=-N$ it becomes
\begin{equation}
 \frac12C_N'(-N)+\sum_{j=0}^{N-1}
 C_j(-N)\zeta(1-2N+2j,a/2)=W\rho_N.
 \label{dg:residue-transport}
\end{equation}
For generic $b,c$, the Gamma term has constant coefficient
\[
 \frac{\rho_N}{2}
 \{\psi(N+1)-\psi(d_0+N)-\psi(b-N)-\psi(c-N)\}.
\]
The resonant Hurwitz term has constant coefficient
$C_N'(-N)/2-\rho_N\psi(a/2)$. Use
\eqref{dg:residue-transport} to combine this derivative with all
lower counterterms. Since $Wd_0=0$,
\[
 W\rho_N=\frac{\rho_N}{2}
 \{\psi(b)-\psi(b-N)+\psi(c)-\psi(c-N)\}.
\]
The shifted digammas cancel and yield
\eqref{dg:antecedent-closed}. Its final arguments are positive, so
normal convergence and holomorphy extend it across all zeros of
$\rho_N$. These antecedent formulas reproduce the needed part of
the earlier resonance construction~\cite{RepoDougall,RepoTransport}.
\end{proof}

\begin{proposition}[Termwise collapse on every integer residue component]
\label{dg:polynomial}
Suppose $a,b,c>0$ and $d_0>0$. For $1\leq k\leq N$, at $c=k$,
\begin{align}
 r_n(-N)&=x_n
  \prod_{r=1}^{k-1}\bigl(x_n^2-(\eta-r)^2\bigr)
  \prod_{r=0}^{N-k-1}
        \bigl(x_n^2-(\eta+1-b+r)^2\bigr).
 \label{dg:grid-factor}
\end{align}
Consequently $F_{N,n}(a,b,k)=0$ for every $n$, and the same statement
holds with $b$ and $c$ interchanged. Locally along any intersection
$b=h$, $c=k$ with $h,k\in\{1,\ldots,N\}$, the holomorphic function
$F_{N,n}$ is divisible by $(b-h)(c-k)$. Its quotient and every fixed
derivative admit locally normally convergent sums. More globally, the
quotient by $(1-b)_N(1-c)_N$ has a holomorphic extension throughout
the stated parameter domain, and its sum is
\begin{align}
 &\sum_{n\geq0}\frac{F_{N,n}(a,b,c)}{(1-b)_N(1-c)_N}\notag\\
 &\quad=\frac{(-1)^N(d_0)_N}{N!}
 \left\{\psi(a/2)+
 \frac{\psi(N+1)-\psi(b)-\psi(c)-\psi(d_0+N)}2\right\}.
 \label{dg:divided-grid}
\end{align}
At a zero of either denominator factor, the summands in this formula
mean their holomorphic extensions just described.
\end{proposition}
\begin{proof}
The Gamma functional equation gives
\[
 r_n(-N)=x_n(n+1)_{k-1}(n+1+a-k)_{k-1}
                 (n+1+a-b)_{N-k}(n+b+k-N)_{N-k}.
\]
Pair the factors symmetric about $x_n$ to obtain
\eqref{dg:grid-factor}. This is an odd polynomial of degree
$2N-1$ in $x_n$. Uniqueness of asymptotic coefficients shows that its
centered coefficients are exactly $C_0(-N),\ldots,C_{N-1}(-N)$,
whereas $C_N(-N)=0$ and every subsequent coefficient is zero.
The identity includes the finitely many small indices where an
individual reciprocal Gamma factor vanishes.

For the divisibility statement, the holomorphic function vanishes
when $b=h$ and when $c=k$. Successive one-variable Taylor expansions
therefore factor $(b-h)(c-k)$. Cauchy's formula on a slightly larger
parameter polydisc transfers the $O(n^{-3})$ uniform remainder bound
to the divided function and to every prescribed derivative. This also
proves normal convergence after division; a quotient cannot be justified
merely by dividing two limiting sums at their common zero.
Every zero of $(1-b)_N$ or $(1-c)_N$ is simple. The other local
factors are units, so this establishes extension of the full divided
sum. Its value follows from \eqref{dg:antecedent-closed} off the
zero loci and then by holomorphy on them.
\end{proof}

The result is stronger than the vanishing of the residue polynomial
or of the summed remainder. In particular, every derivative tangent
to one of these components still has zero remainder termwise.
Mixed transverse derivatives are the first potentially nonzero ones.
For example, write
\[
 B_N(a,b,c)=\psi(a/2)+
 \frac{\psi(N+1)-\psi(b)-\psi(c)-\psi(d_0+N)}2.
\]
At $b=h$, $c=k$ with $1\leq h,k\leq N$,
\begin{align}
 \sum_{n\geq0}\partial_b\partial_cF_{N,n}(a,h,k)
 ={}&\frac{(-1)^{N+h+k}(d_0)_N}{N!}
 (h-1)!(N-h)!(k-1)!(N-k)!\notag\\
 &\times B_N(a,h,k),\qquad d_0=1+a-h-k.
 \label{dg:grid-first-mixed}
\end{align}
This uses $\left.\partial_b(1-b)_N\right|_{b=h}
=(-1)^h(h-1)!(N-h)!$ and keeps all finite-head terms.

\subsection{A two-parameter generating identity for all positive shifts}

Write $b=1+p$, $c=1+q$, $A=a-1$, and keep $u=-1$. Put
\begin{align}
 Q_n(a;p,q)
 &=\frac{\Gamma(n+a)\Gamma(n+1+p)\Gamma(n+1+q)
                         \Gamma(n+a-p-q)}
 {\Gamma(n+1)\Gamma(n+a-p)\Gamma(n+a-q)
                         \Gamma(n+1+p+q)},\label{dg:Q}\\
 B_a(p,q)&=\psi(a/2)+\frac{\psi(2)-\psi(1+p)-\psi(1+q)
                                      -\psi(a-p-q)}2.
 \label{dg:B}
\end{align}
Thus $Q_n(a;p,0)=Q_n(a;0,q)=1$, and $r_n(-1)=x_nQ_n(a;p,q)$.

\begin{theorem}[Double-zero generating function]
\label{dg:generator}
For every $a>0$ and sufficiently small complex $p,q$,
\begin{align}
 \sum_{n=0}^{\infty}\left\{
 x_n\bigl(Q_n(a;p,q)-1\bigr)
       +\frac{(A-p-q)pq}{x_n}\right\}
 =-(A-p-q)pq\,B_a(p,q).
 \label{dg:generator-eq}
\end{align}
The series is locally normally convergent in $(a,p,q)$ in a
neighborhood of each point $(a,0,0)$ with $\Re a>0$. It may be
differentiated at every fixed mixed order. Both sides are divisible
by $pq$, and the divided identity also holds normally at $p=q=0$.
\end{theorem}
\begin{proof}
For $a>1$ and small $p,q$, Lemma~\ref{dg:antecedent} applies with
$N=1$. Its residue coefficient is
\[
 C_1(-1)=-(a-1-p-q)pq,
\]
and the closed resonance value is exactly the right side of
\eqref{dg:generator-eq}. The centered expansion supplies
\[
 Q_n(a;p,q)=1-\frac{(A-p-q)pq}{x_n^2}+O(n^{-4})
\]
locally uniformly, so the displayed bracket is $O(n^{-3})$.
All Gamma arguments in \eqref{dg:Q} lie in a pole-free
neighborhood when $\Re a>0$ and $p,q$ are sufficiently small; the
same centered expansion is uniform on compact subsets there.
Thus the sum and the right side are holomorphic across $a=1$ and
throughout this connected positive-half-plane germ. The identity
theorem extends their equality from $a>1$ to all $a>0$.
For any prescribed compact $a$-set one can choose a smaller common
polydisc in $p,q$ and apply Cauchy's coefficient formula to the
uniform remainder. This proves the derivative and division claims.
\end{proof}

The restriction $d_0>0$ was needed only in the original region from
which continuation was proved. After the exact cancellation above,
$a=1$ is a regular point, not an excluded parameter or a formal limit.

\subsection{A finite Bell formula at every mixed order}

For $k\geq2$, define
\[
 \Lambda_k(n;a)=(-1)^k\psi^{(k-1)}(n+a)
                                -\psi^{(k-1)}(n+1).
\]
Taylor expansion of the four Gamma ratios gives the exact germ
\begin{equation}
 \log Q_n(a;p,q)=\sum_{r,s\geq1}
       \Lambda_{r+s}(n;a)\frac{p^rq^s}{r!s!}.
 \label{dg:logQ}
\end{equation}
Define a finite bivariate Bell polynomial by
\[
 \mathcal B_{rs}(n;a)=r!s![p^rq^s]
 \exp\left(\sum_{i,j\geq1}
       \Lambda_{i+j}(n;a)\frac{p^iq^j}{i!j!}\right).
\]
Only $i\leq r$, $j\leq s$, and powers at most $\min(r,s)$
of the exponent are needed. In particular,
\begin{align*}
 \mathcal B_{11}&=\Lambda_2,&
 \mathcal B_{21}&=\Lambda_3,&
 \mathcal B_{31}&=\Lambda_4,\\
 \mathcal B_{22}&=\Lambda_4+2\Lambda_2^2,&
 \mathcal B_{32}&=\Lambda_5+6\Lambda_2\Lambda_3,\\
 \mathcal B_{33}&=\Lambda_6+9\Lambda_2\Lambda_4
                            +9\Lambda_3^2+6\Lambda_2^3.
\end{align*}

Set $B_{rs}=\partial_p^r\partial_q^sB_a(0,0)$, with
\[
 B_{00}=\psi(a/2)+\frac{1+\gamma-\psi(a)}2,
\]
and, for $r+s=k\geq1$,
\begin{equation}
 B_{rs}=-\frac12\left\{
 \mathbf1_{s=0}\psi^{(r)}(1)+
 \mathbf1_{r=0}\psi^{(s)}(1)+(-1)^k\psi^{(k)}(a)\right\}.
 \label{dg:B-jets}
\end{equation}
Terms carrying a negative subscript below are omitted.

\begin{corollary}[Complete mixed parameter table]
\label{dg:all-parameter}
For $r,s\geq1$, the single bracketed series
\begin{align}
 S_{rs}(a)=\sum_{n\geq0}\left\{
 x_n\mathcal B_{rs}(n;a)
 +\frac{A\mathbf1_{r=s=1}
  -2\mathbf1_{(r,s)=(2,1)}-2\mathbf1_{(r,s)=(1,2)}}{x_n}
 \right\}
 \label{dg:S}
\end{align}
converges absolutely, and
\begin{equation}
 S_{rs}(a)=-rs\left\{
 A B_{r-1,s-1}-(r-1)B_{r-2,s-1}-(s-1)B_{r-1,s-2}
 \right\}.
 \label{dg:all-parameter-eq}
\end{equation}
For every $r+s\geq4$, no counterterm is needed: the original
polygamma polynomial $x_n\mathcal B_{rs}(n;a)$ is summable by itself.
\end{corollary}
\begin{proof}
Differentiate \eqref{dg:generator-eq} and use
\eqref{dg:logQ}. The counterterm is a polynomial of total
degree three in $p,q$, which explains its exact disappearance at
total derivative order four. Leibniz's rule gives
\eqref{dg:all-parameter-eq}; normal convergence was proved
before differentiation, including the small indices.
\end{proof}

For example, at $a=2$ the first rows are
\[
\begin{array}{c|c|c}
 (r,s)&\mathcal B_{rs}&S_{rs}(2)\\\hline
 (1,1)&\Lambda_2&0\\
 (2,1)&\Lambda_3&1\\
 (3,1)&\Lambda_4&-6\zeta(3)\\
 (2,2)&\Lambda_4+2\Lambda_2^2&-4\zeta(3)\\
 (3,2)&\Lambda_5+6\Lambda_2\Lambda_3&24\zeta(3)-18\zeta(4)\\
 (3,3)&\Lambda_6+9\Lambda_2\Lambda_4+9\Lambda_3^2+6\Lambda_2^3
                                    &108(\zeta(4)-\zeta(5)).
\end{array}
\]
The first two rows use the counterterms in \eqref{dg:S}.
Every later row is an ordinary convergent sum without subtraction.

\subsection{A closed square of trigamma differences}

\begin{theorem}[Weighted trigamma square]
\label{dg:square}
For every real $a>0$,
\begin{align}
 &\sum_{n=0}^{\infty}\left(n+\frac a2\right)
       \bigl(\psi'(n+a)-\psi'(n+1)\bigr)^2\notag\\
 &\qquad=\frac{a-1}{4}\psi''(a)
     +\frac{\psi'(a)-\zeta(2)}2
     +\frac{3(a-1)}2\zeta(3).
 \label{dg:square-eq}
\end{align}
The series converges locally normally for $\Re a>0$, with an
ordinary algebraic square in that complex extension.
\end{theorem}
\begin{proof}
The $(2,2)$ and $(3,1)$ rows give respectively
\begin{align*}
 \sum_nx_n(\Lambda_4+2\Lambda_2^2)
   &=2A\psi''(a)+4\psi'(a)-4\zeta(2),\\
 \sum_nx_n\Lambda_4
   &=\frac{3A}{2}\{\psi''(1)+\psi''(a)\}
                                  +3\{\psi'(a)-\zeta(2)\}.
\end{align*}
Subtract and divide by two, using $\psi''(1)=-2\zeta(3)$.
Every series used in the subtraction is absolutely convergent.
Independently, the polygamma asymptotic or its convergent series
shows $\Lambda_2(n;a)=O(n^{-2})$ uniformly on compact subsets of
$\Re a>0$. Thus the weighted square is $O(n^{-3})$; this also
directly verifies the domain and regularity at $a=1$.
\end{proof}

At $a=3/2$, the functional equation and duplication formula~\cite{DLMFPolygamma} give
\[
 \psi'(3/2)=3\zeta(2)-4,\qquad
 \psi''(3/2)=-14\zeta(3)+16.
\]
The right side of \eqref{dg:square-eq} becomes
\[
 \frac18(-14\zeta(3)+16)
 +\frac12(2\zeta(2)-4)+\frac34\zeta(3)
 =\zeta(2)-\zeta(3).
\]
For $H_n^{(r)}=\sum_{k=1}^nk^{-r}$, $H_0^{(r)}=0$,
\[
 \psi'(n+1)-\psi'(n+3/2)
 =2\{2H_{2n+1}^{(2)}-H_n^{(2)}-\zeta(2)\}.
\]
Therefore
\begin{equation}
 \boxed{\sum_{n=0}^{\infty}(4n+3)
 \{2H_{2n+1}^{(2)}-H_n^{(2)}-\zeta(2)\}^{2}
 =\zeta(2)-\zeta(3).}
 \label{dg:half-harmonic-square}
\end{equation}
No regularized interpretation is involved in this positive series.

The continuation to $a=1/2$ gives a companion identity,
\begin{equation}
 \sum_{n=0}^{\infty}(4n+1)
 \{\zeta(2)-2H_{2n}^{(2)}+H_n^{(2)}\}^{2}
 =\zeta(2)+\zeta(3).
 \label{dg:other-half-harmonic-square}
\end{equation}
Here $\psi'(1/2)=3\zeta(2)$ and
$\psi''(1/2)=-14\zeta(3)$. At $a=2$, the square formula reduces to
$\sum_{k\geq1}k^{-3}=\zeta(3)$, a useful elementary normalization.

The $(3,3)$ row also supplies a higher harmonic example. At $a=2$,
writing $k=n+1$,
\[
 \Lambda_2=-k^{-2},\quad \Lambda_4=-6k^{-4},\quad
 \Lambda_6=-120k^{-6},\quad
 \Lambda_3=2\{2(\zeta(3)-H_{k-1}^{(3)})-k^{-3}\}.
\]
Substitution gives
\begin{equation}
 \boxed{\sum_{k=1}^{\infty}k
 \{2(\zeta(3)-H_{k-1}^{(3)})-k^{-3}\}^2
 =3\zeta(4)-\zeta(5).}
 \label{dg:cubic-tail-square}
\end{equation}

\subsection{Every resonant polynomial moment from the same germ}

For $N\geq1$, $s=p+q$, set
\[
 P_{N-1}(x;s)=\prod_{r=0}^{N-2}\{x^2-(a/2-s+r)^2\},\qquad P_0=1.
\]
The exact Gamma functional equation gives
\begin{equation}
 r_n(-N)=x_nP_{N-1}(x_n;p+q)Q_n(a;p,q).
 \label{dg:N-transport}
\end{equation}
It retains every finite-head zero as a factor before derivatives are
taken. This identity and the two-variable Bell germ generate the
summands at every resonance without singular logarithmic derivatives.

More explicitly, write
\[
 P_{N-1}(x;s)=\sum_{k=0}^{N-1}P_k(s)x^{2N-2-2k},\qquad
 Q_n(a;p,q)\sim\sum_{j\geq0}D_j(p,q)x_n^{-2j}.
\]
Here $D_0=1$ and every $D_j$, $j\geq1$, is divisible by $pq$.
The counterterms are the finite convolution
\begin{equation}
 C_j(-N)=\sum_{k=0}^{\min(j,N-1)}P_k(p+q)D_{j-k}(p,q).
 \label{dg:N-coefficients}
\end{equation}
Equation~\eqref{dg:antecedent-closed} becomes
\begin{align}
 \sum_{n\geq0}F_{N,n}(a,1+p,1+q)
 &=\frac{(-1)^N}{N!}(a-1-p-q)_N(-p)_N(-q)_N\notag\\
 &\quad\times B_N(a,1+p,1+q).
 \label{dg:N-all}
\end{align}
This holds as a normally convergent germ for every $a>0$, by the
same continuation argument as Theorem~\ref{dg:generator}.
Taking any prescribed $[p^rq^s]$ is a finite exact algorithm involving
only the coefficients in \eqref{dg:N-coefficients}, polygamma
values, and finite generalized harmonic numbers from the Pochhammer
factors. Formula \eqref{dg:N-all} is a representation theorem;
it does not assert that all possible products of polygamma functions
can be separated individually by these mixed derivatives.

\subsection{A mixed spectral germ containing every Stieltjes order}

We now keep one derivative in each parameter direction and permit
the spectral parameter to move: $u=-1+t$. For $a>0$, let
\begin{align}
 R_n(a;t)&=x_n\frac{\Gamma(n+1)\Gamma(n+a-t)}
                       {\Gamma(n+a)\Gamma(n+1+t)},\notag\\
 L_n(a;t)&=\psi(n+1)-\psi(n+a-t)+\psi(n+a)-\psi(n+1+t),\notag\\
 K_n(a;t)&=\psi'(n+a-t)-\psi'(n+1+t).
 \label{dg:spectral-summands}
\end{align}
At $p=q=0$, these satisfy
\[
 \partial_p\partial_q r_n(-1+t)
                 =R_n(a;t)\{L_n(a;t)^2+K_n(a;t)\}.
\]
The identity uses the unspecialized summand before derivatives.

To state the endpoint without a division by $a-1$, define
\begin{align}
 F_a(t)&=\frac{\Gamma(a-t)}{\Gamma(a)\Gamma(1+t)(1-t)},\notag\\
 V_a(t)&=\psi(1)-\psi(a-t)+\psi(a)-\psi(1+t),\notag\\
 W_a(t)&=\psi'(a-t)-\psi'(a),\notag\\
 D_a(t)&=F_a(t)\left\{2t(1+tV_a(t))
 -(a-1)\bigl[(1+tV_a(t))^2+t^2W_a(t)\bigr]\right\}.
 \label{dg:D}
\end{align}
All of these are holomorphic near $t=0$ for every $a>0$, including
$a=1$, and $D_a(0)=-(a-1)$.

\begin{theorem}[Complete mixed spectral--Stieltjes germ]
\label{dg:spectral}
For $a>0$ and sufficiently small complex $t$, the series
\begin{equation}
 M_a(t)=\sum_{n=0}^{\infty}\left\{
 R_n(a;t)\bigl(L_n(a;t)^2+K_n(a;t)\bigr)
 +(a-1-2t)x_n^{-1-2t}\right\}
 \label{dg:M}
\end{equation}
converges locally normally and equals
\begin{equation}
 M_a(t)=\frac{D_a(t)}{2t}
             +(a-1-2t)\zeta(1+2t,a/2).
 \label{dg:M-eq}
\end{equation}
The apparent singularity at $t=0$ is removable. Define $\gamma_m(v)$
by the Laurent expansion at the simple Hurwitz pole~\cite{DLMFHurwitz}:
\[
 \zeta(1+z,v)=z^{-1}+\sum_{m\geq0}\frac{(-1)^m}{m!}\gamma_m(v)z^m.
\]
Then every fixed derivative has the explicit expression
\begin{equation}
 M_a^{(m)}(0)=\frac{D_a^{(m+1)}(0)}{2(m+1)}
 -\mathbf1_{m=0}
 +(-2)^m\{(a-1)\gamma_m(a/2)+m\gamma_{m-1}(a/2)\}.
 \label{dg:all-Stieltjes}
\end{equation}
The $m\gamma_{m-1}$ term is omitted when $m=0$. Every derivative
on the left can be taken in the single convergent bracket
\eqref{dg:M}.
\end{theorem}
\begin{proof}
For $a>1$, differentiate \eqref{dg:antecedent-master} with $K=2$
once in $p$ and once in $q$. The lower counterterm $C_0=1$ is
parameter independent. Direct differentiation of the cubic Bernoulli
coefficient gives
\[
 \left.\partial_p\partial_qC_1(-1+t)\right|_{p=q=0}=-(a-1)+2t.
\]
The Gamma term is $E_1(t;p,q)/(2t)$, where
\[
 E_1(t;p,q)=
 \frac{\Gamma(1+p)\Gamma(1+q)\Gamma(a-p-q-t)\Gamma(1+t)}
 {\Gamma(a-1-p-q)\Gamma(1+p+t)\Gamma(1+q+t)}
                  \frac{(p+t)(q+t)}{t-1}.
\]
The factors $(p+t)(q+t)$ must be retained: setting $p=q=0$
before differentiating would erase the mixed derivative.
Twice applying the product rule and using
\[
 \Gamma(a)=(a-1)\Gamma(a-1),\quad
 \psi(a-1)=\psi(a)-(a-1)^{-1},\quad
 \psi'(a-1)=\psi'(a)+(a-1)^{-2}
\]
shows that $\partial_p\partial_qE_1(t;0,0)=D_a(t)$.
The resulting expression \eqref{dg:D} contains no singular
factor at $a=1$. This proves \eqref{dg:M-eq} for $a>1$.
Uniform Stirling remainders on a slightly enlarged parameter
neighborhood give a bound
$O(n^{-3-2\Re t}(1+\log^j n))$ for every fixed $t$-derivative of
the bracket. The displayed functions are holomorphic for $\Re a>0$
near $t=0$, so continuation gives all $a>0$ as before.

Finally subtract the pole explicitly:
\[
 M_a(t)=\frac{D_a(t)+(a-1)-2t}{2t}
 +(a-1-2t)\sum_{h\geq0}\frac{(-2)^h}{h!}\gamma_h(a/2)t^h.
\]
Taylor coefficient extraction yields
\eqref{dg:all-Stieltjes}, including the constant $-1$ at
$m=0$. The formula is for derivatives rather than divided Taylor
coefficients; its factor $1/(m+1)$ is therefore essential.
\end{proof}

For an explicit first spectral derivative, write
\[
 \ell=1+\gamma-\psi(a),\quad
 d=\psi'(a)-\zeta(2),\quad
 \mu_n=-\psi(n+a)-\psi(n+1).
\]
Then \eqref{dg:all-Stieltjes} gives the ordinary identity
\begin{align}
 &\sum_{n=0}^{\infty}\left\{
 x_n\bigl[\Lambda_3(n;a)+\mu_n\Lambda_2(n;a)\bigr]
       -\frac{2+2(a-1)\log x_n}{x_n}\right\}\notag\\
 &\quad=2\psi(a/2)+\ell
       -\frac{a-1}{4}\bigl(\ell^2+5d+1\bigr)
       -2(a-1)\gamma_1(a/2).
 \label{dg:first-spectral}
\end{align}
Indeed the coefficient of $t^2$ in $D_a(t)$ is
$2\ell-(a-1)(\ell^2+5d+1)/2$.
At $a=2$, the right side reduces to
$1-\gamma^2-2\gamma_1$. Writing $k=n+1$, this is equivalently
\begin{align}
 \sum_{k=1}^{\infty}\left\{
 4k(\zeta(3)-H_{k-1}^{(3)})-\frac1{k^2}-\frac2k
 +\frac{2(H_{k-1}-\gamma-\log k)}k\right\}
 =1-\gamma^2-2\gamma_1.
 \label{dg:harmonic-Stieltjes}
\end{align}
The divergent-looking pieces in this last formula are one bracket;
the proof supplies absolute convergence of their combination.

\subsection{Further questions suggested by the mixed normal form}

The arbitrary-order formulas above solve the requested first
two-parameter table at $b=c=1$ and provide its spectral extension.
Several sharper questions remain. First, at a fixed homogeneous
order, which individual polygamma-product moments are recoverable
from the finitely many mixed Bell combinations? Their number need
not match the number of all possible products, so the all-order
generating theorem alone does not prove complete separation.
Second, the polynomial formula on $b=h$, $c=k$ supplies starting
points at every residue-grid intersection; compact explicit tables
away from $(1,1)$ should retain their finite-head zeros and avoid
singular logarithmic derivative notation. Third, at rational $a$,
the Stieltjes germ can be expressed in finite residue-class or
Dirichlet-character coordinates; it is natural to determine which
sectors cancel in particular mixed orders. Finally, primitive or
generating-variable versions of the mixed identities can be tested
for reductions to colored polylogarithms, with the function class
and endpoint normalization stated before any claimed reduction.

% END INLINED SOURCE: sections/03_dougall.tex
\clearpage
% BEGIN INLINED SOURCE: sections/04_twists.tex
\section{Unequal twists: an exact normal form and its finite reductions}
\label{tw:sec}

The incoming twisted-Stieltjes report asks for a normal form when two
frequency shifts are genuinely different
\cite[research question on unequal twists]{RepoTwisted}.
An equal-twist Gamma quotient cannot by itself answer that question.
We give three complementary answers: a convergent expansion in ordinary
Lerch families, an all-index convolution formula with one bounded
correction kernel, and a precise finite-reduction criterion for the
underlying scalar functions. The last criterion concerns functional
identities; it makes no inference about independence of sampled periods.

\subsection{Fourier conventions and the entire family}

Write $\T=\R/\Z$, $e_n(x)=e^{2\pi\ii nx}$, and
$\widehat U(n)=\langle U,e_{-n}\rangle$. Let $\theta,\eta$ be real
nonintegers, with $\theta\ne\eta$, and set
\begin{align}
 \lambda_{n,\theta}&=2\pi\ii(n+\theta),\\
 L_{n,\theta}&=\log(2\pi|n+\theta|)
                +\frac{\pi\ii}{2}\sgn(n+\theta).
 \label{tw:log}
\end{align}
Powers of $\lambda_{n,\theta}$ use this logarithm. The normalized
single-twist family and the mixed family are defined by
\begin{align}
 \widehat W_s^\theta(n)&=\lambda_{n,\theta}^{s-1},\label{tw:W}\\
 \widehat{\mathcal K_{\alpha,\beta}^{\theta,\eta}}(n)
   &=\lambda_{n,\theta}^{\alpha}\lambda_{n,\eta}^{\beta}.
 \label{tw:K}
\end{align}
Both are entire with values in periodic distributions, jointly in all
displayed spectral variables. Indeed their Fourier coefficients and every
fixed number of order derivatives have polynomial growth, uniformly on
compact spectral sets. Pairing with a smooth test function therefore gives
a normally convergent series. In particular,
\begin{equation}
 \mathcal K_{\alpha,\beta}^{\theta,\eta}
   =W_{1+\alpha}^{\theta}*W_{1+\beta}^{\eta}.
 \label{tw:Kconv}
\end{equation}
For $\Re(\alpha+\beta)<-1$, the Fourier series itself is absolutely
and uniformly convergent. The distributional definition supplies all
other spectral orders without assigning an ordinary value to a singular
endpoint.

For $0<\theta<1$, the standard Lerch family gives, on $0<x<1$,
\begin{equation}
 W_s^\theta(x)=\frac{e^{-2\pi\ii\theta x}}{\Gamma(1-s)}
       \Phi(e^{-2\pi\ii\theta},s,x).
 \label{tw:lerch}
\end{equation}
This identity is understood through analytic continuation when necessary.
It follows initially by unfolding the Fourier coefficient of
$\sum_{m\ge0}e^{-(\varepsilon+2\pi\ii\theta)(m+x)}(m+x)^{-s}$
and then letting $\varepsilon\downarrow0$. This is also the
normalization established in the incoming report. The Lerch function,
its relation $\Li_s(z)=z\Phi(z,s,1)$, and its analytic continuation
are classical \cite{DLMFLerch}. Other real twists follow by integer
modulation, $W_s^{\theta+k}=e_{-k}W_s^\theta$.

\subsection{A convergent Lerch expansion for arbitrary orders}

Put $d=\theta-\eta$ and $h=2\pi\ii d$. Choose a finite set
$F\subset\Z$ such that, for some $\rho<1$,
\begin{equation}
 \left|\frac{d}{n+\eta}\right|\le\rho\qquad(n\notin F).
 \label{tw:Fcondition}
\end{equation}
For example, all modes with $|n+\eta|\le2|d|$ may be put in $F$.
Let $\Pi_FU=\sum_{n\in F}\widehat U(n)e_n$.

\begin{theorem}[Normal convergence in single-twist Lerch coordinates]
\label{tw:tailtheorem}
For all $\alpha,\beta\in\C$,
\begin{align}
 \mathcal K_{\alpha,\beta}^{\theta,\eta}
 ={}&\sum_{n\in F}
        \lambda_{n,\theta}^{\alpha}\lambda_{n,\eta}^{\beta}e_n
       \notag\\
 &+\sum_{k=0}^{\infty}\binom\alpha k h^k
       (I-\Pi_F)W_{1+\alpha+\beta-k}^{\eta}.
 \label{tw:tailseries}
\end{align}
The series converges normally after pairing with any smooth test
function, locally uniformly in $(\alpha,\beta)$. Every fixed order
derivative in either spectral variable may be taken termwise. After
finitely many initial terms are removed, it converges normally in any
prescribed $C^q$ norm, locally on spectral parameter sets.
\end{theorem}

\begin{proof}
For $n\notin F$, condition \eqref{tw:Fcondition} implies that
$n+\theta$ and $n+\eta$ have the same sign. Therefore there is no
branch correction in
\[
 L_{n,\theta}=L_{n,\eta}
             +\log\left(1+\frac{d}{n+\eta}\right).
\]
The ordinary binomial theorem proves the equality of every Fourier
coefficient in \eqref{tw:tailseries}. On a compact set of $\alpha$,
$|\binom\alpha k|\le C(1+k)^M$ for suitable $C,M$. This elementary
bound follows either from the finite product or by Cauchy's formula
applied to $(1+z)^\alpha$ on a disc of radius between $\rho$ and
$1$; the latter also gives a geometric bound sufficient here.
Each summand is consequently bounded by a polynomial in $|n|$
times a summable geometric sequence in $k$. The test-function Fourier
coefficients absorb the polynomial. Cauchy's formula on a slightly
larger compact spectral set justifies all order derivatives.

For the $C^q$ assertion, take $k_0$ larger than
$q+1+\sup\Re(\alpha+\beta)$. Factor
$|\lambda_{n,\eta}|^{\Re(\alpha+\beta)-k_0}$ from the terms
with $k\ge k_0$. The remaining ratios are at most a fixed geometric
sequence, while the $q$ derivatives in $x$ still leave an absolutely
summable series over $n$. This proves the stated normal convergence.
\end{proof}

The finitely many removed modes are essential when the twists lie in
different integer intervals. Omitting them can change both the phase and
the value of the result. The theorem gives a convergently evaluable
normal form at every spectral order, including all Lerch order jets.

\subsection{One bounded kernel controls every mixed Stieltjes jet}

Let $P_\theta$ have Fourier coefficient
$\gamma+L_{n,\theta}$, and define
\begin{equation}
 B_\theta(u)=\Gamma(1-u)W_{1+u}^{\theta}
   =\delta_0+\sum_{m\ge0}\frac{(-1)^{m+1}}{m!}
                         U_m^\theta u^{m+1}.
 \label{tw:B}
\end{equation}
Thus $P_\theta=-U_0^\theta$. This agrees with the incoming
finite-part Lerch normalization. Introduce the difference
\begin{equation}
 A_{\theta,\eta}=P_\theta-P_\eta,
 \qquad\widehat A_{\theta,\eta}(n)=L_{n,\theta}-L_{n,\eta}.
 \label{tw:A}
\end{equation}

\begin{lemma}[The difference is an ordinary bounded kernel]
\label{tw:bounded}
If $\theta,\eta$ belong to the same open interval between consecutive
integers, then
\begin{equation}
 A_{\theta,\eta}(x)=2\pi\ii\int_\eta^\theta
       \frac{e^{-2\pi\ii t x}}{1-e^{-2\pi\ii t}}\dd t,
       \qquad 0<x<1.
 \label{tw:Aintegral}
\end{equation}
This function extends smoothly to each closed side of the endpoint,
is bounded, and has jump
\begin{equation}
 A_{\theta,\eta}(0+)-A_{\theta,\eta}(1-)
                    =2\pi\ii(\theta-\eta).
 \label{tw:Ajump}
\end{equation}
For arbitrary real noninteger endpoints, the same formula holds with
the path indented below the intervening integers. Equivalently, for
$\eta<\theta$, it is the Cauchy principal value of the real integral
plus $\pi\ii\sum_{k\in\Z\cap(\eta,\theta)}e_{-k}(x)$.
\end{lemma}

\begin{proof}
The Green kernel
\begin{equation}
 G_t(x)=\frac{e^{-2\pi\ii t x}}{1-e^{-2\pi\ii t}}
 \label{tw:G}
\end{equation}
has Fourier coefficient $1/(2\pi\ii(n+t))$. Integrating that
coefficient gives the logarithmic difference in \eqref{tw:A} if
no zero is crossed. Smoothness in $x$ up to each side follows by
differentiating the compact $t$-integral. Subtracting the endpoint
values gives \eqref{tw:Ajump}.

When $t$ crosses $-n$, continuation below the pole adds $+\pi\ii$
to the logarithmic difference, precisely the principal-logarithm phase
change from $-\pi/2$ to $+\pi/2$ in \eqref{tw:log}.
The residue of $2\pi\ii G_t(x)$ at an integer $k$ is $e_{-k}(x)$.
Subtracting those finitely many simple poles defines an ordinary bounded
function plus a finite Fourier polynomial. Its Fourier coefficients
are again exactly \eqref{tw:A}. Reversing orientation covers
$\theta<\eta$.
\end{proof}

For $0<\theta,\eta<1$, an alternative useful expression is
\begin{equation}
 A_{\theta,\eta}(x)
 =e^{-2\pi\ii\eta x}\Phi(e^{-2\pi\ii\eta},1,x)
  -e^{-2\pi\ii\theta x}\Phi(e^{-2\pi\ii\theta},1,x).
 \label{tw:ALerch}
\end{equation}
The two $1/x$ singularities cancel. If $\theta=p/q\in(0,1)$,
the corresponding term reduces to ordinary rational-shift digammas:
\begin{equation}
 e^{-2\pi\ii px/q}\Phi(e^{-2\pi\ii p/q},1,x)
 =-\frac{e^{-2\pi\ii px/q}}q
      \sum_{j=0}^{q-1}e^{-2\pi\ii pj/q}
                          \psi\left(\frac{x+j}{q}\right).
 \label{tw:rational}
\end{equation}
To prove it, split the Lerch defining sum into residue classes:
\[
 \Phi(e^{-2\pi\ii p/q},s,x)
   =q^{-s}\sum_{j=0}^{q-1}e^{-2\pi\ii pj/q}
                          \zeta\left(s,\frac{x+j}{q}\right).
\]
Continue this identity to $s=1$. The Hurwitz poles cancel because the finite character sum is
zero; $\gamma_0(a)=-\psi(a)$ gives the formula.

Write $\exp_*(uA)=\delta_0+\sum_{r\ge1}u^r A^{*r}/r!$.
Since $A$ is bounded and belongs to $L^1(\T)$, this exponential is
entire with values in $\delta_0+L^1(\T)$, by the usual convolution
norm estimate. Put
\begin{equation}
 C(u,v)=\frac{\Gamma(1-u)\Gamma(1-v)}{\Gamma(1-u-v)}.
 \label{tw:C}
\end{equation}

\begin{theorem}[All-index unequal-twist formula]
\label{tw:jets}
As analytic germs of periodic distributions,
\begin{equation}
 B_\theta(u)*B_\eta(v)
       =C(u,v)B_\eta(u+v)*\exp_*(uA_{\theta,\eta}).
 \label{tw:mixedB}
\end{equation}
Consequently every mixed finite-part product is the finite coefficient
formula
\begin{equation}
 U_m^\theta*U_n^\eta
  =(-1)^{m+n}m!n!\,[u^{m+1}v^{n+1}]
       C(u,v)B_\eta(u+v)*\exp_*(uA_{\theta,\eta}).
 \label{tw:alljets}
\end{equation}
Only powers $A^{*r}$ with $0\le r\le m+1$ occur. The $r=0$
part is the known equal-twist Gamma table; the other terms have Fourier
coefficients tending to zero and introduce no new point-supported
counterterm.
\end{theorem}

\begin{proof}
At a Fourier mode, the right side of \eqref{tw:mixedB} is
\[
 \Gamma(1-u)\Gamma(1-v)
      e^{(u+v)L_{n,\eta}}
      e^{u(L_{n,\theta}-L_{n,\eta})},
\]
which is exactly the left side. This proves the analytic identity by
normal convergence after testing. Taking the indicated coefficients
proves \eqref{tw:alljets}. For $r\ge1$, the Fourier coefficient of
$A^{*r}$ is $O(|n|^{-r})$, whereas a fixed spectral jet of
$B_\eta$ grows only as a fixed polynomial in $\log|n|$. Thus their
product tends to zero. A nonzero distribution supported at one point
is a finite sum of derivatives of its Dirac mass and has polynomial
Fourier coefficients; none can be added while preserving these
coefficients. This statement refers to the displayed canonical
decomposition, not to an arbitrary splitting of a distribution into
summands.
\end{proof}

The first symmetric identity is particularly compact:
\begin{equation}
 \boxed{\quad
 U_0^\theta*U_0^\eta
   =U_1^\theta+U_1^\eta-\zeta(2)\delta_0
                     -\frac12 A_{\theta,\eta}^{*2}.
 \quad}
 \label{tw:firstmixed}
\end{equation}
Indeed $2U_1^\theta=P_\theta^{*2}+\zeta(2)\delta_0$,
and polarization gives the result. Here $A^{*2}$ is an ordinary
continuous function, since its Fourier coefficients are $O(n^{-2})$.
The new term is therefore an explicitly evaluable ordinary convolution
of bounded Lerch differences, while the contact coefficient remains
fixed. An equal-twist formula with the correction omitted is false:
at mode $n$ it misses $\tfrac12(L_{n,\theta}-L_{n,\eta})^2$.

\subsection{Finite reductions and their precise boundary}

\begin{proposition}[Integer resolvent products]
\label{tw:partial}
For positive integers $p,q$, with $h=2\pi\ii(\theta-\eta)$,
\begin{align}
 \mathcal K_{-p,-q}^{\theta,\eta}
 ={}&\sum_{j=1}^p
  \frac{(-1)^q\binom{p+q-j-1}{p-j}}{h^{p+q-j}}
                  W_{1-j}^\theta\notag\\
 &+\sum_{j=1}^q
  \frac{(-1)^{q-j}\binom{p+q-j-1}{q-j}}{h^{p+q-j}}
                  W_{1-j}^\eta.
 \label{tw:partialformula}
\end{align}
If one of $\alpha,\beta$ is a nonnegative integer, there is likewise
a finite binomial reduction at the other twist. All integer pairs
therefore reduce to finitely many single-twist families. In particular,
\begin{equation}
 G_\theta*G_\eta=\frac{G_\eta-G_\theta}{2\pi\ii(\theta-\eta)}.
 \label{tw:resolvent}
\end{equation}
\end{proposition}

\begin{proof}
Expand $(A-h)^{-q}$ at $A=0$ and $(B+h)^{-p}$ at $B=0$, with
$A-B=h$. The displayed coefficients are exactly the principal parts
of $1/(A^pB^q)$ at its two poles. Their difference from that rational
function is an entire rational function vanishing at infinity and is
therefore zero. Substitute $A=\lambda_{n,\theta}$ and
$B=\lambda_{n,\eta}$. The polynomial cases follow by the finite
binomial theorem. Equation \eqref{tw:resolvent} is $p=q=1$.
\end{proof}

For positive integer $j$, the terms in \eqref{tw:partialformula}
are elementary on the interval:
\begin{equation}
 W_{1-j}^\theta(x)
  =\frac{(-1)^{j-1}}{(j-1)!(2\pi\ii)^{j-1}}
      \partial_\theta^{j-1}
      \frac{e^{-2\pi\ii\theta x}}{1-e^{-2\pi\ii\theta}}.
 \label{tw:elementary}
\end{equation}
The combined right side has its usual confluent limit
$W_{1-p-q}^{\theta}$ as $\eta\to\theta$, although individual
partial-fraction coefficients diverge. For example,
\begin{align}
 G_{1/4}*G_{3/4}(x)
  &=\frac{\ii}{2\pi}
    \left((1+\ii)e^{-3\pi\ii x/2}
                   -(1-\ii)e^{-\pi\ii x/2}\right),\label{tw:quarters}\\
 G_{1/2}^{*2}(x)&=e^{-\pi\ii x}
                         \left(\frac{x}{2}-\frac14\right).
 \label{tw:half}
\end{align}

\begin{theorem}[Finite scalar reduction criterion]
\label{tw:classification}
Fix distinct complex numbers $a,b$ and coherent branches on a simply
connected domain avoiding $-a,-b$. Consider
$f(z)=(z+a)^\alpha(z+b)^\beta$. It is a finite linear combination
of single-centre powers and order jets
\[
 (z+a)^r\Log^k(z+a),\qquad (z+b)^t\Log^\ell(z+b),
 \qquad k,\ell\in\N\cup\{0\},
\]
with arbitrary complex exponents and constant coefficients, if and only if
\begin{equation}
 \alpha\in\Z_{\ge0},\qquad\text{or}\qquad
 \beta\in\Z_{\ge0},\qquad\text{or}\qquad
 \alpha,\beta\in\Z.
 \label{tw:criterion}
\end{equation}
\end{theorem}

\begin{proof}
Sufficiency follows from the binomial theorem and rational partial
fractions. For necessity, suppose first $\alpha\notin\Z$ and
$\beta\notin\Z_{\ge0}$. Near $z=-a$, write $w=z+a$ and
$c=b-a\ne0$. Then
\[
 f(z)=w^\alpha\sum_{j\ge0}\binom\beta j c^{\beta-j}w^j.
\]
All infinitely many coefficients are nonzero. Terms based at $-b$
are holomorphic near $-a$. Apply analytic continuation once around
$w=0$ and subtract the original expression. On the left this
multiplies the displayed infinite series by
$e^{2\pi\ii\alpha}-1\ne0$; on the proposed right it removes every
term holomorphic there and leaves only finitely many powers times
polynomials in $\Log w$. Let $E=w\,d/dw$. A nonzero polynomial $P(E)$ annihilates that
finite expression: for each occurring $w^r\Log^k w$, include the
factor $(E-r)^{k+1}$. Applied to the left, it gives a nonzero constant
multiple of
\[
 w^\alpha\sum_{j\ge0}P(\alpha+j)
                    \binom\beta j c^{\beta-j}w^j.
\]
Vanishing would force $P(\alpha+j)=0$ for every $j\ge0$, contrary
to $P$ being a nonzero polynomial. This is a contradiction.

Interchanging the two centres proves the same conclusion if
$\beta\notin\Z$ and $\alpha\notin\Z_{\ge0}$. Unless
\eqref{tw:criterion} holds, at least one of these two cases occurs.
\end{proof}

The qualification ``scalar'' matters. Equality on the discrete set of
Fourier modes is weaker than equality of the underlying holomorphic
functions; the theorem does not promote the former to the latter.
It explains exactly when the elementary functional mechanism behind
\eqref{tw:partialformula} can terminate. Arbitrary orders still have
the normally convergent expansion \eqref{tw:tailseries}.

\subsection{Argument derivatives and canonical primitives}

Let $D_\theta=d/dx+2\pi\ii\theta$. Fourier multiplication gives
the exact contiguous and differential laws
\begin{align}
 D_\theta\mathcal K_{\alpha,\beta}^{\theta,\eta}
   &=\mathcal K_{\alpha+1,\beta}^{\theta,\eta},&
 D_\eta\mathcal K_{\alpha,\beta}^{\theta,\eta}
   &=\mathcal K_{\alpha,\beta+1}^{\theta,\eta},\label{tw:diff}\\
 \mathcal K_{\alpha+1,\beta}^{\theta,\eta}
  -\mathcal K_{\alpha,\beta+1}^{\theta,\eta}
   &=h\mathcal K_{\alpha,\beta}^{\theta,\eta}.
 \label{tw:contiguous}
\end{align}
Since neither twist is integral, $D_\theta$ and $D_\eta$ are
invertible on periodic distributions. Their inverses lower the
corresponding spectral exponent. These statements remain valid after
any order derivatives in $\alpha,\beta$.

For an ordinary integrable $f$, the unique periodic covariant primitive
has the explicit form
\begin{equation}
 (D_\eta^{-1}f)(x)=e^{-2\pi\ii\eta x}
 \left\{\int_0^x e^{2\pi\ii\eta t}f(t)\dd t
   +\frac{e^{-2\pi\ii\eta}}{1-e^{-2\pi\ii\eta}}
          \int_0^1e^{2\pi\ii\eta t}f(t)\dd t\right\}.
 \label{tw:primitive}
\end{equation}
Differentiation proves $D_\eta U=f$, and the displayed constant is
forced by $U(1)=U(0)$. This supplies ordinary quadratures for every
bounded correction in \eqref{tw:alljets}, while the singular
single-twist part retains its canonical distributional primitive.


% END INLINED SOURCE: sections/04_twists.tex
\clearpage
% BEGIN INLINED SOURCE: sections/05_collisions.tex
\section{Simultaneous Stieltjes collisions and their finite anomalies}
\label{cl:sec}

The incoming separated-triple theorem assumes pairwise disjoint singular
grids. Its first further question asks whether simultaneous collision
constants agree with a directly merged ambient finite part
\cite[Question~1]{RepoExact}. The bilinear coincidence report supplies
such a comparison for two factors \cite{RepoCoincident}. The triple
problem contains an additional phenomenon: two singular factors can
multiply a resonant Taylor coefficient of a third regular factor.
We prove an explicit finite discrepancy, identify its dependence on the
collision rates, and then give an all-factor local completion theorem.
The distinction between cutoff and spectral prescriptions also appears
in the different endpoint setting of \cite{RepoEndpoint}; here the
collision parameter moves singularities in the real argument.

\subsection{Fixed coordinates and the geometric collision theorem}
\label{cl:geometric}

Put $f(x)=\gamma_0(x)=-\psi(x)$ and $\gamma=\gamma_0(1)$. For
$0<a<b$ and $0<\epsilon<1/b$, let
\[
 C_{a,b}(\epsilon)=\operatorname{FP}_x\int_0^1
 f(x)f(\{x+a\epsilon\})f(\{x+b\epsilon\})\,dx.
\]
At each of the three distinct grid points $0$, $1-b\epsilon$, and
$1-a\epsilon$, the finite part deletes a right interval of ambient
length $\eta$ and takes the coefficient independent of $\eta$.
The three cutoff variables can also be independent. Define the merged
moment in the same coordinate by
\begin{equation}\label{cl:merged}
 Q=\operatorname{FP}_x\int_0^1f(x)^3\,dx.
\end{equation}
This has the ordinary convergent representation
\begin{equation}\label{cl:merged-convergent}
 Q=\int_0^1\left[f(x)^3-\frac1{x^3}-\frac{3\gamma}{x^2}
       -\frac{3(\gamma^2-\zeta(2))}{x}\right]dx-\frac12-3\gamma.
\end{equation}

\begin{theorem}[Triple collision constant, with its rate dependence]
\label{cl:collision}
With $L=\log\epsilon$, define
\begin{align*}
 P_{-2}(a,b;L)&=\frac{L}{ab}
    +\frac{\log a}{a(b-a)}-\frac{\log b}{b(b-a)},\\
 P_{-1}(a,b;L)&=\gamma\left[
    \frac{L+\log a}{a}+\frac{L+\log b}{b}
    +\frac{L+\log(b-a)}{b-a}\right],\\
 P_0(a,b;L)&=\zeta(2)\left[
    \frac{a-b}{a}(L+\log a)
    +\frac{b-a}{b}(L+\log b)
    +\frac b{b-a}(L+\log(b-a))\right].
\end{align*}
Then
\begin{equation}\label{cl:rate-expansion}
 C_{a,b}(\epsilon)=\epsilon^{-2}P_{-2}(a,b;L)
    +\epsilon^{-1}P_{-1}(a,b;L)+P_0(a,b;L)+Q
    +O(\epsilon|\log\epsilon|).
\end{equation}
Consequently its coefficient independent of $\epsilon$ and $\log\epsilon$
is $Q+\zeta(2)D(a,b)$, where
\[
 D(a,b)=\frac{a-b}{a}\log a+\frac{b-a}{b}\log b
               +\frac b{b-a}\log(b-a).
\]
In particular,
\begin{align}
 C_{1,2}(\epsilon)
 ={}&\frac{L-\log2}{2\epsilon^2}
   +\frac{\gamma(5L+\log2)}{2\epsilon}
   +\frac{\zeta(2)}2(3L+\log2)+Q
   +O(\epsilon|\log\epsilon|),\label{cl:equally-spaced}\\
 \operatorname{CT}_{\epsilon}C_{1,2}(\epsilon)
 ={}&Q+\frac{\zeta(2)}2\log2\ne Q.\label{cl:finite-anomaly}
\end{align}
Here $\operatorname{CT}_{\epsilon}$ removes divergent powers and all
nonconstant logarithmic terms; it does not remove finite constants
inside a chosen collision counterterm. Removing the complete $P_0$,
including its finite part, does recover $Q$.
\end{theorem}

\begin{proof}
Choose a fixed small $\delta>0$ and work in the circle coordinate
$x\in(-\delta,\delta)$ around the merged point. The Hurwitz shift
identity \cite[Equation~25.11.3]{DLMFHurwitz} gives the exact local representation
\[
 f(\{x\})=\frac{\mathbf 1_{x>0}}x+h(x),\qquad
 h(x)=f(1+x)=\gamma-\zeta(2)x+x^2R(x),
\]
where $R$ is analytic near zero. For the three rates $(0,a,b)$ write
$S_c(x)=\mathbf 1_{x+c\epsilon>0}/(x+c\epsilon)$ and
$H_c(x)=h(x+c\epsilon)$. Expand the product of the three $S_c+H_c$.
The all-regular term and the three terms with exactly one $S_c$ are
analytic in $\epsilon$ after the prescribed one-sided finite parts
are taken. Their limits are their merged counterparts. The integral
outside this neighborhood is also analytic for small $\epsilon$.

For a term with two singular factors, let their rates be $c<d$ and
let the remaining regular factor have rate $e$. Set
$t=x+c\epsilon$, $\Delta=d-c$, $\rho=e-c$ and $\ell=\delta+c\epsilon$.
Its local integral is
\[
 \operatorname{FP}\int_0^\ell
     \frac{h(t+\rho\epsilon)}{t(t+\Delta\epsilon)}\,dt.
\]
The two necessary elementary primitives are
\begin{align*}
 A_0&=\operatorname{FP}\int_0^\ell\frac{dt}{t(t+\Delta\epsilon)}
 =\frac{\log\ell-\log(\ell+\Delta\epsilon)
                  +\log(\Delta\epsilon)}{\Delta\epsilon},\\
 A_1&=\int_0^\ell\frac{dt}{t+\Delta\epsilon}
 =\log(\ell+\Delta\epsilon)-\log(\Delta\epsilon).
\end{align*}
The constant and linear parts of $h$ contribute
$\gamma A_0-\zeta(2)(A_1+\rho\epsilon A_0)$. As $\epsilon\downarrow0$,
this equals
\[
 \frac{\gamma\log(\Delta\epsilon)}{\Delta\epsilon}
 +\zeta(2)\frac{d-e}{d-c}\log(\Delta\epsilon)
 -\frac\gamma\delta-\zeta(2)\log\delta+O(\epsilon|\log\epsilon|).
\]
The last two displayed constants are exactly the merged finite parts
of $\gamma x^{-2}-\zeta(2)x^{-1}$ on $(0,\delta)$.
The analytic quadratic remainder converges to its merged ordinary
integral. More precisely, use
\[
 \frac{(t+\rho\epsilon)^2}{t(t+\Delta\epsilon)}
 =1+\frac{(2\rho-\Delta)\epsilon}{t+\Delta\epsilon}
       +\frac{\rho^2\epsilon^2}{t(t+\Delta\epsilon)}
\]
and the analytic, locally bounded function $R$ to obtain an error
$O(\epsilon|\log\epsilon|)$. Thus the three pairs give exactly
$\epsilon^{-1}P_{-1}+P_0$ beyond their merged values.

The three-singular term is supported on $x>0$. Partial fractions give
\[
 \frac1{x(x+a\epsilon)(x+b\epsilon)}
 =\frac1{ab\epsilon^2x}
  -\frac1{a(b-a)\epsilon^2(x+a\epsilon)}
  +\frac1{b(b-a)\epsilon^2(x+b\epsilon)}.
\]
Its lower-end contribution is exactly $\epsilon^{-2}P_{-2}$. The
upper-end expression tends to $-1/(2\delta^2)$, which is precisely
$\operatorname{FP}\int_0^\delta x^{-3}\,dx$; its remaining error is
$O(\epsilon)$. Summing all eight terms and the nonsingular exterior
integral proves the formula. Only simple logarithms occur because the
zero-index singular pieces are rational functions with simple linear
factors; no $\log^2\epsilon$ term has been suppressed.
\end{proof}

The rate dependence can be isolated without an unevaluated merged
constant:
\begin{equation}\label{cl:rate-difference}
 \operatorname{CT}_{\epsilon}C_{1,3}(\epsilon)
 -\operatorname{CT}_{\epsilon}C_{1,2}(\epsilon)
 =\zeta(2)\left(\log2+\frac23\log3\right).
\end{equation}
These are positive shifts, exactly as in the definition of $C_{a,b}$.
For three digamma factors all displayed correlations change sign.

\subsection{Resonant subsets and a holomorphic completion}
\label{cl:subsets}

This section supplies an all-index explanation of the new phenomenon.
Put
\[
 G(z,x)=\zeta(1-z,x)+\frac1z
     =\sum_{m\ge0}\gamma_m(x)\frac{z^m}{m!},\qquad
 c_r(z)=(-1)^r(1-z)_r,
\]
and write
\[
 g_r(z,x)=\partial_x^rG(z,x),\qquad h_r(z)=g_r(z,1).
\]
In particular $h_0(z)=\zeta(1-z)+1/z$ and
$h_r(z)=c_r(z)\zeta(r+1-z)$ for $r\ge1$. All these functions are
holomorphic near $z=0$.

Let there be $d\ge1$ factors, positive integer frequencies $p_i$, real
shifts $a_i$, and nonnegative derivative orders $r_i$. Let
$\mathcal B$ be the finite union of their singular grids on the circle,
and define the collision set at $b\in\mathcal B$ by
\[
 C_b=\{i:p_i b+a_i\in\mathbb Z\}.
\]
For $t=x-b>0$ put
\[
 H_{b,j}(z_j,t)=
 \begin{cases}
 g_{r_j}(z_j,1+p_jt),&j\in C_b,\\
 g_{r_j}(z_j,\{p_jb+a_j\}+p_jt),&j\notin C_b.
 \end{cases}
\]
For every nonempty subset $S\subseteq C_b$ set
\[
 \Lambda_S=\sum_{i\in S}z_i,\qquad
 k_S=\sum_{i\in S}(r_i+1)-1,
\]
and define the complete resonant numerator
\[
 A_{b,S}(\boldsymbol z)=
 \left(\prod_{i\in S}c_{r_i}(z_i)p_i^{z_i-r_i-1}\right)
 [t^{k_S}]\prod_{j\notin S}H_{b,j}(z_j,t).
\]
Empty products are one. For $d\ge2$, the term with
$S=\{1,\dots,d\}$ vanishes because $k_S>0$.

\begin{theorem}[All-factor local completion]
\label{cl:allfactor}
Let $I_{\boldsymbol r}(\boldsymbol z)$ be the meromorphic continuation
of the ordinary integral
\[
 \int_0^1\prod_{i=1}^d
      g_{r_i}(z_i,\{p_i x+a_i\})\,dx,
\]
initially defined, for example, by $\Re z_i>r_i+1$ for every $i$.
Then
\begin{equation}\label{cl:completion}
 J_{\boldsymbol r}(\boldsymbol z)
 =I_{\boldsymbol r}(\boldsymbol z)
       -\sum_{b\in\mathcal B}\ \sum_{\varnothing\ne S\subseteq C_b}
              \frac{A_{b,S}(\boldsymbol z)}{\Lambda_S}
\end{equation}
is holomorphic near the origin. Its coefficient
\[
 \left(\prod_i m_i!\right)[z_1^{m_1}\cdots z_d^{m_d}]
              J_{\boldsymbol r}(\boldsymbol z)
\]
is exactly the ambient-coordinate finite part of the corresponding
product of $\gamma_{m_i}^{(r_i)}$. The theorem permits arbitrary partial
and complete collisions.
\end{theorem}

\begin{proof}
The exact Hurwitz shift gives each local factor as
\[
 g_{r_i}(z_i,\{p_ix+a_i\})
 =c_{r_i}(z_i)p_i^{z_i-r_i-1}t^{z_i-r_i-1}
       +H_{b,i}(z_i,t)\quad(i\in C_b).
\]
Expand the finite product by its nonempty subsets of singular pieces.
For $S$, expand its regular complement through degree $k_S$. The
Taylor coefficient of degree $j$ integrates meromorphically with
denominator $\Lambda_S+j-k_S$. All $j<k_S$ have denominator nonzero
near the origin and already generate their coordinate finite parts.
At $j=k_S$ the term is $A_{b,S}t^{\Lambda_S-1}$.
Its meromorphic integral on $(0,\ell)$ is
$A_{b,S}\ell^{\Lambda_S}/\Lambda_S$, while its coefficientwise
coordinate finite-part generator is
$A_{b,S}(\ell^{\Lambda_S}-1)/\Lambda_S$. Their difference is the
complete numerator $A_{b,S}/\Lambda_S$, not a residue with its
parameter dependence deleted. The Taylor remainder has exponent
$\Lambda_S$ and is integrable uniformly on a sufficiently small
spectral polydisc; all fixed spectral derivatives add only logarithmic
powers. Finite summation over grid points proves holomorphy and the
coefficient interpretation, and also constructs the local meromorphic
continuation uniquely from the initial convergent integral.
\end{proof}

For two coincident factors, a product of both singular pieces has no
regular complement of the positive degree needed for resonance. Hence
the only polar divisors through the origin are $z_i=0$. For three
coincident factors, pairs of singular pieces have one regular
complement, so $z_i+z_j=0$ contributes for the first time. This is the
structural difference behind the triple collision anomaly.

\subsection{The explicit same-point triple generator}
\label{cl:triple}

\begin{corollary}[Every coincident triple Stieltjes jet]
\label{cl:triple-completion}
At unit frequencies and a common endpoint, put $R=r_1+r_2+r_3$.
The complete generator is
\begin{align}
 J_{\boldsymbol r}(\boldsymbol z)=I_{\boldsymbol r}(\boldsymbol z)
 &-\sum_{i=1}^3\frac{c_{r_i}(z_i)}{z_i}
     \sum_{u+v=r_i}\frac{h_{r_j+u}(z_j)h_{r_k+v}(z_k)}{u!v!}\notag\\
 &-\sum_{i<j}\frac{c_{r_i}(z_i)c_{r_j}(z_j)h_{R+1}(z_k)}
          {(r_i+r_j+1)!(z_i+z_j)}.\label{cl:triple-generator}
\end{align}
In each summand $\{i,j,k\}=\{1,2,3\}$.
For $r_i=0$ this is particularly short:
\begin{equation}\label{cl:triple-zero}
 J_{\boldsymbol0}=I_{\boldsymbol0}
 -\sum_i\frac{h_0(z_j)h_0(z_k)}{z_i}
 -\sum_{i<j}\frac{h_1(z_k)}{z_i+z_j}.
\end{equation}
\end{corollary}
\begin{proof}
For a singleton $S=\{i\}$, its resonant Taylor degree is $r_i$,
and Leibniz's formula gives the inner sum in the first line.
For a pair $S=\{i,j\}$, the resonant degree is $r_i+r_j+1$.
The remaining factor has derivative order $r_k$, so its Taylor
coefficient is $h_{R+1}(z_k)/(r_i+r_j+1)!$. The three-singular
subset has no regular complement and makes no resonant contribution.
Theorem~\ref{cl:allfactor} now gives both formulas.
\end{proof}
The last sum is indispensable: $h_1(0)=-\zeta(2)\ne0$.

\subsubsection{A completely specified Tornheim realization}
\label{cl:tornheim}

The classical Tornheim function is initially
\[
 T(A,B,C)=\sum_{m,n\ge1}\frac1{m^A n^B(m+n)^C}.
\]
Let $\alpha_i=z_i-r_i$ and define
\begin{align*}
 K_3(\boldsymbol\alpha)
 &=\frac{2\prod_i\Gamma(\alpha_i)}{(2\pi)^{\sum_i\alpha_i}}
   \sum_{k=1}^3
   \cos\!\left(\frac\pi2(\alpha_i+\alpha_j-\alpha_k)\right)
        T(\alpha_i,\alpha_j,\alpha_k),\\
 K_2(A,B)&=\frac{2\Gamma(A)\Gamma(B)}{(2\pi)^{A+B}}
          \cos\!\left(\frac\pi2(A-B)\right)\zeta(A+B).
\end{align*}
The classical Fourier calculation \cite{EspinosaMoll} gives these as the integrals of,
respectively, three or two factors $\zeta(1-\alpha,x)$ in an initial
absolute-convergence domain. Its six sign cones explain the three
cosines. For the regularized Stieltjes generators one has the exact
meromorphic identity
\[
 I_{\boldsymbol r}(\boldsymbol z)
 =\left(\prod_i c_{r_i}(z_i)\right)K_3(\boldsymbol\alpha)
  +\sum_{k:r_k=0}\frac{c_{r_i}(z_i)c_{r_j}(z_j)}{z_k}
                  K_2(\alpha_i,\alpha_j)
  +\frac{\mathbf 1_{\boldsymbol r=\boldsymbol0}}{z_1z_2z_3}.
\]
Terms with two scalar $1/z_i$ factors vanish because the remaining
Hurwitz factor has mean zero in the initial domain. This combined
formula, followed by the complete local subtraction above, is a
finite all-index coefficient rule. No arithmetic reduction of the
Tornheim jets to single-zeta coordinates is asserted.

Meromorphic continuation of $T$ near the required nonpositive integer
orders can be proved directly from
\[
 T(A,B,C)=\frac1{\Gamma(C)}\int_0^\infty
 t^{C-1}\operatorname{Li}_A(e^{-t})\operatorname{Li}_B(e^{-t})\,dt.
\]
On a small interval insert the classical local expansion \cite[Equation~25.12.12]{DLMFPolylog}
\[
 \operatorname{Li}_A(e^{-t})
 =\Gamma(1-A)t^{A-1}
       +\sum_{j\ge0}\zeta(A-j)\frac{(-t)^j}{j!},
\]
which is uniform on compact parameter neighborhoods avoiding positive
integer $A$. Finite subtraction and termwise integration give explicit
linear polar divisors; the remaining Mellin integral is holomorphic.
The same argument handles the colored kernels with one or both twists
equal to one. Thus an entire-order claim valid for two nontrivial
twists must be weakened to meromorphic continuation when collisions
force trivial twists. The incoming separated-grid theorem is valid
on its stated hypotheses; it is its unrestricted extrapolation that
would be incorrect.

\subsection{Spectral rays and Laurent expansion chambers}
\label{cl:spectral}

\begin{proposition}[Exact cubic spectral-ray correction]
\label{cl:rays}
Take $r_i=0$ and $z_i=c_i t$, where $c_i\ne0$ and $c_i+c_j\ne0$.
Since the completed $J$ is holomorphic, its constant is always $Q$.
The uncompleted meromorphic integral instead has constant
\begin{align}
 \operatorname{CT}_{t}I_{\boldsymbol0}(c_1t,c_2t,c_3t)
 ={}&Q+\gamma\gamma_1\sum_i\frac{c_j+c_k}{c_i}\notag\\
       &+(\zeta(2)+\zeta'(2))\sum_{i<j}\frac{c_k}{c_i+c_j}.
       \label{cl:ray-constant}
\end{align}
Here $\gamma_1$ is the ordinary first Stieltjes constant and $\zeta'$
is the spectral derivative. On the diagonal the discrepancy is
\[
 6\gamma\gamma_1+\frac32\bigl(\zeta(2)+\zeta'(2)\bigr).
\]
\end{proposition}
\begin{proof}
Use \eqref{cl:triple-zero} and the expansions
$h_0(z)=\gamma+\gamma_1z+O(z^2)$ and
$h_1(z)=-\zeta(2)+(\zeta(2)+\zeta'(2))z+O(z^2)$.
After division by $c_it$ or $(c_i+c_j)t$, precisely the linear
numerator coefficient contributes to the constant. The holomorphic
remainder contributes $J_{\boldsymbol0}(0,0,0)=Q$. For the diagonal,
set $c_1=c_2=c_3=1$.
\end{proof}
This operation concerns spectral rays and is distinct from the
geometric collision $\epsilon\to0$ proved earlier. Mixed poles also
mean that a rectangular Laurent expansion is no longer a canonical
operation unless an expansion chamber or an iterated prescription is
specified. The holomorphic completion avoids either choice.

\begin{proposition}[Three chamber constants for one derivative moment]
\label{cl:chambers}
For $\boldsymbol r=(1,0,0)$ one has
\begin{equation}\label{cl:derivative-moment}
 J_{1,0,0}(0,0,0)
 =\operatorname{FP}\int_0^1\gamma_0'(x)\gamma_0(x)^2\,dx
 =2\gamma\zeta(2)-\zeta(3).
\end{equation}
In an iterated Laurent
chamber ordering the three moduli, the raw constant is instead
\begin{equation}\label{cl:chamber-values}
 \operatorname{CT}_{\rm chamber}I_{1,0,0}
 =\begin{cases}
 \zeta(3),& |z_1|\text{ is largest},\\
 0,& |z_1|\text{ is intermediate},\\
 -\zeta(3),& |z_1|\text{ is smallest}.
 \end{cases}
\end{equation}
On the spectral diagonal the same raw integral is identically zero:
$I_{1,0,0}(t,t,t)=0$ as a meromorphic function.
\end{proposition}
\begin{proof}
The local expansion is
\[
 \gamma_0(x)=x^{-1}+\gamma-\zeta(2)x+\zeta(3)x^2+O(x^3).
\]
Consequently the coefficient of $x^0$ in $\gamma_0(x)^3$ is
$\gamma^3-6\gamma\zeta(2)+3\zeta(3)$. Integrate the exact derivative
$(\gamma_0^3)'/3$ from the cutoff to $1$ and retain its constant.
This proves \eqref{cl:derivative-moment} directly, without a spectral
regularization.

For clarity, write $I_{1,0,0}-J_{1,0,0}=S+P$, separating singleton
and pair terms in \eqref{cl:triple-generator}. Since $c_1(z_1)=z_1-1$,
these are
\begin{align*}
 S={}&\frac{z_1-1}{z_1}
       \bigl[h_1(z_2)h_0(z_3)+h_0(z_2)h_1(z_3)\bigr]\\
   &+\frac{h_1(z_1)h_0(z_3)}{z_2}
     +\frac{h_1(z_1)h_0(z_2)}{z_3},\\
 P={}&\frac{(z_1-1)h_2(z_3)}{2(z_1+z_2)}
     +\frac{(z_1-1)h_2(z_2)}{2(z_1+z_3)}
     +\frac{h_2(z_1)}{z_2+z_3}.
\end{align*}
The singleton constant is $-2\gamma\zeta(2)$ in every chamber.
Using $h_2(0)=2\zeta(3)$, the only pair terms contributing to a
constant are
\[
 \zeta(3)\frac{z_1}{z_1+z_2}
       +\zeta(3)\frac{z_1}{z_1+z_3}.
\]
In a chamber with $|z_j|<|z_1|$, expand
$z_1/(z_1+z_j)=\sum_{k\ge0}(-z_j/z_1)^k$; its constant is one.
In the reversed chamber its expansion starts with $z_1/z_j$ and
has constant zero. Thus $P$ contributes $N\zeta(3)$, where
$N$ is the number of other variables dominated by $z_1$.
Adding the holomorphic constant in \eqref{cl:derivative-moment} gives
$(N-1)\zeta(3)$, proving all three cases.

Finally, for $\Re t>2$, the function $G(t,x)$ has equal endpoint
values and a continuous first derivative on the period. Therefore
\[
 I_{1,0,0}(t,t,t)=\int_0^1
       \partial_xG(t,x)G(t,x)^2\,dx
       =\frac{G(t,1)^3-G(t,0)^3}{3}=0.
\]
Meromorphic continuation preserves this identity. In particular,
its spectral diagonal constant is zero even though its coordinate
finite part is \eqref{cl:derivative-moment}.
\end{proof}

In contrast, for $\boldsymbol r=\boldsymbol0$ and unit frequencies,
the polar numerators do not depend on the variables in their own
denominator. Their nonnegative rectangular coefficients vanish in
every standard chamber. Thus the general mixed-pole warning should
not be misread as asserting chamber dependence of the raw constant
in that special zero-derivative case; its spectral-ray constants still
differ, as shown above.

\subsection{An exact cancellation at every number of factors}

The incoming question on odd derivative order identifies equal Stieltjes
indices and one total argument derivative as its first target
\cite[Question~3]{RepoExact}. That case has a complete local answer.
The mechanism is the elementary fundamental theorem of calculus, applied
before taking the cutoff constant; keeping its endpoint term gives the
following more general identity.

\begin{theorem}[The all-factor total-derivative projection]
\label{cl:total-derivative}
Let $d\ge1$, $m_1,\ldots,m_d\ge0$, and put
$Z_0=\{i:m_i=0\}$ and $g_i(t)=\gamma_{m_i}(1+t)$. Then
\begin{align}
 &\FP\int_0^1\sum_{i=1}^d\gamma_{m_i}'(x)
                   \prod_{j\ne i}\gamma_{m_j}(x)\,dx\notag\\
 &\qquad=-\sum_{\varnothing\ne S\subseteq Z_0}
              [t^{|S|}]\prod_{j\notin S}g_j(t).
 \label{cl:total-projection}
\end{align}
The sum on the right is empty when all $m_i>0$. Each surviving
coefficient is a finite expression in ordinary Stieltjes constants and
argument derivatives of Stieltjes functions at $1$.
\end{theorem}
\begin{proof}
For every fixed cutoff, the integrand is the ordinary derivative of
$\prod_i\gamma_{m_i}(x)$. Hence its finite-part integral is
\[
 \prod_i\gamma_{m_i}(1)
       -\CT_{x\downarrow0}\prod_i\gamma_{m_i}(x).
\]
Use the exact shift
$\gamma_{m_i}(x)=x^{-1}\log^{m_i}x+g_i(x)$ and expand by the
subset $S$ of singular pieces. The empty subset contributes exactly
$\prod_i\gamma_{m_i}(1)$, canceling the upper endpoint.
A nonempty subset contributes
$x^{-|S|}\log^{\sum_{i\in S}m_i}x\prod_{j\notin S}g_j(x)$.
It has a constant in both $x$ and $\log x$ precisely from the
Taylor degree $|S|$ and only when $S\subseteq Z_0$.
This proves \eqref{cl:total-projection}. All individual finite parts
exist by the same finite local power--log expansion, so taking the
constant after the finite sum introduces no interchange issue.
\end{proof}

For equal positive indices the theorem gives the requested cancellation
at every number of factors:
\begin{equation}\label{cl:equal-index-cancellation}
 \boxed{\quad
 \FP\int_0^1\gamma_m'(x)\gamma_m(x)^{d-1}\,dx=0,
       \qquad d\ge1,\quad m\ge1.
 \quad}
\end{equation}
For the zero index, with $h(t)=\gamma_0(1+t)$ as above, the complete
remaining local value is
\begin{equation}\label{cl:zero-index-projection}
 \FP\int_0^1\gamma_0'(x)\gamma_0(x)^{d-1}\,dx
 =-\frac1d\sum_{k=1}^d\binom dk[t^k]h(t)^{d-k}.
\end{equation}
This is a finite zeta polynomial because
$h(t)=\gamma+\sum_{j\ge1}(-1)^j\zeta(j+1)t^j$ near zero.
For $d=3$ it recovers \eqref{cl:derivative-moment}; for $d=4$ it gives
\[
 \FP\int_0^1\gamma_0'(x)\gamma_0(x)^3\,dx
 =3\gamma^2\zeta(2)-3\gamma\zeta(3)-\frac{11}{4}\zeta(4).
\]
Thus the symmetric projection removes the multiple spectral terms
entirely in this one-derivative case, leaving its explicitly specified
local endpoint contribution. Higher odd derivative projections remain
a separate question.

\subsection{Repository scope and independent numerical diagnostics}
\label{cl:numerics}

The inspected snapshot is
\texttt{dcd95baeb7c0e914b138bddef6829c5b1d52b726}.
The incoming \src{Exact_Identities} report, section
\texttt{triple\_dilation.tex}, assumes pairwise disjoint grids throughout.
Its \texttt{audit\_research.tex}, Question~1, asks to remove that
restriction and compare simultaneous collision constants with merged
finite parts. The subset theorem completes its fixed-collision local
subtraction question, and the explicit $(0,\epsilon,2\epsilon)$
calculation shows that its suggested matching is not valid for plain
geometric constant extraction. A finite rate-dependent correction
is necessary. The preceding bilinear coincidence results are not
contradicted.

The supplied script \src{code/verify_collisions.py} computes the
separated moments by three independent convergent interval integrals.
For an interval of length $\ell$ starting at the singular point of one
factor, with the other two arguments $\eta_1+t,\eta_2+t$, it uses
\[
 \int_0^\ell\left[
 \frac{f(\eta_1+t)f(\eta_2+t)-f(\eta_1)f(\eta_2)}{t}
 +h(t)f(\eta_1+t)f(\eta_2+t)\right]dt
 +f(\eta_1)f(\eta_2)\log\ell.
\]
This is exactly the ambient finite part with the simple pole visibly
removed. It does not implement the collision expansion or the
Tornheim representation. The merged moment is evaluated by a local
digamma Laurent series plus ordinary quadrature, with two independent
split points. At 70-digit working precision the computed values are
\[
 Q=-1.9662627343915343406643753384116314238258689764917\ldots,
\]
and
\[
 \frac12\zeta(2)\log2=0.5700907053214263787372899294961746726083149206823\ldots.
\]
Numerical residuals tend to zero at the proved
$O(\epsilon|\log\epsilon|)$ rate for $(a,b)=(1,2),(1,3),(2,5)$.
For the equally spaced example, the residual means
$C_{1,2}(\epsilon)-\epsilon^{-2}P_{-2}
 -\epsilon^{-1}P_{-1}-P_0-Q$; the complete finite term in $P_0$
is included in this diagnostic:
\begin{center}
\begin{tabular}{@{}rr@{}}
\toprule
$\epsilon$ & Numerical residual\\
\midrule
$10^{-6}$ & $-9.332630963\times10^{-5}$\\
$10^{-8}$ & $-1.237723998\times10^{-6}$\\
$10^{-10}$ & $-1.542186195\times10^{-8}$\\
\bottomrule
\end{tabular}
\end{center}
Omitting the finite term would instead leave a limit
$\zeta(2)\log2/2$. The full record contains seventeen collision
diagnostics, the two-split merged calculation, working precision,
and the finite correction. It is reproduced from the package root by
\begin{verbatim}
python code/verify_collisions.py --output results/collisions.json
\end{verbatim}
These are floating-point diagnostics, not certified interval enclosures;
the proof is the exact local decomposition above.

The classical inputs are the Hurwitz shift, argument-derivative and
Fourier formulas \cite{DLMFHurwitz}, the local polylogarithm expansion
\cite{DLMFPolylog}, and the Hurwitz--Tornheim product relation
\cite{EspinosaMoll}. The contributions here are the complete
coordinate-normalized subtraction, the finite geometric collision
anomaly, and their exact comparison with the specified spectral
prescriptions. The statements do not claim a reduction of every
Tornheim derivative to ordinary zeta values.

% END INLINED SOURCE: sections/05_collisions.tex
\clearpage
% BEGIN INLINED SOURCE: sections/06_audit_research.tex
\section{Source corrections, verification, and further research}
\label{audit:sec}

\subsection{Corrections proposed against the pinned manuscript}

The inspection was targeted: the relevant harmonic, Dougall, twist,
collision, and Gaussian-status material was checked, rather than every
chapter and every archived computational receipt. The following two
edits are proposed for \src{chapters/03-algebraic.tex} in the
canonical manuscript \cite{RepoMain}. They concern the exact snapshot
specified in Section~\ref{scope:sec}. The package contains a small
unified patch; the repository itself has not been changed.

\paragraph{A real logarithm on the negative branch.}
In the equation labeled \src{cleo:eq:lintanh}, put
$\theta=\arcsin\lambda$ for real $|\lambda|<1$. The formula valid
on both sides of zero is
\begin{equation}\label{audit:real-atanh}
 \int_0^{\pi/2}\operatorname{arctanh}(\lambda\sin x)\,dx
 =\theta\log\left|\tan\frac\theta2\right|
       +2\operatorname{Cl}_2(\theta)
       -\frac12\operatorname{Cl}_2(2\theta).
\end{equation}
Here $\operatorname{Cl}_2(\theta)=\sum_{n\ge1}\sin(n\theta)/n^2$.
The value at $\theta=0$ is understood by continuity. The original
display omits the absolute value. For negative $\lambda$, the
real logarithm of $\tan(\theta/2)$ is not defined, whereas the
integral is real. Taking a principal complex logarithm would add
an unwanted imaginary term. This is a carry-forward correction
already identified in the incoming audit material, retained here
because the pinned canonical display still needs it.

For a direct verification, differentiate the left side after
$\lambda=\sin\theta$ to obtain $\theta/\sin\theta$.
The derivative identity
$\operatorname{Cl}_2'(\theta)=-\log|2\sin(\theta/2)|$
shows that the three logarithmic terms in the derivative of the
right side cancel, leaving the same $\theta/\sin\theta$.
Both sides tend to zero at the origin. This proves
\eqref{audit:real-atanh} on the two intervals and through their
continuous joining point.

\paragraph{Remove an unsupported non-elementarity assertion.}
The same chapter describes the $(1,1)$ diagonal as the
``simplest non-elementary diagonal.'' Its exact evaluation
\[
 \int_0^{\pi/2}\arctan^2(\sin x)\,dx
       =\pi\chi_2(3-2\sqrt2)
\]
where $\chi_2(z)=\sum_{n\ge0}z^{2n+1}/(2n+1)^2$,
does not prove non-elementarity or minimality in any specified
constant field. The chapter's later paragraph correctly says that
these properties remain unproved. We propose the neutral wording
``For the $(1,1)$ diagonal, the formula gives \dots''.
The value is unchanged. A finite unsuccessful relation search
cannot justify a stronger assertion.

\subsection{A corrected collision expectation, not a failed source theorem}

The separated-grid theorem in \cite{RepoExact} assumes disjoint
singular supports and is not contradicted by this continuation.
Nor is the proved bilinear coincidence comparison invalidated.
What fails is an unqualified extension of the positive matching
expectation in its Question~1 to geometric triple collisions.
The exact counterexample is \eqref{cl:finite-anomaly}.
A corrected research statement is:

\begin{quote}
For a specified collision path and ambient cutoff, determine the
complete local collision expression, including its finite
rate-dependent term. Subtract that expression before comparing
the limit with the merged ambient finite part.
\end{quote}

For the zero-index triple with arbitrary fixed rates $0<a<b$,
Theorem~\ref{cl:collision} proves precisely this statement, with
the complete expression
$\epsilon^{-2}P_{-2}+\epsilon^{-1}P_{-1}+P_0$.
The all-factor theorem fixes the local meromorphic completion at
a prescribed collision configuration; it does not by itself
compute every geometric collision expansion for arbitrary
Stieltjes indices and rates depending on different powers of
$\epsilon$. Those extensions are questions below.

\subsection{Statements that must retain their current status}

The source's $S_2$ and $S_4$ reductions already have proofs.
The current $S_6$ candidate is labeled
\src{cycloquot:conj:S6} in
\src{chapters/04-cyclotomic-quotients.tex}; the retained $S_8$
candidate is labeled \src{s8new:conj:S8} in
\src{chapters/04-S8-candidate.tex} \cite{RepoMain}.
Both remain conjectural here. The later $S_8$ vector must not
be confused with the different older vector that the manuscript
rigorously rejected. Exact equivalence of two candidate coordinate
vectors proves equivalence, not vanishing of either residual.
Likewise, an interval containing zero proves proximity only.
No new Gaussian integer relation was fitted in this work.

The harmonic Newton theorem solves the explicit-representation
problem for the entire remainder. It leaves the finite-basis
reduction question open at general depth and spectral order.
The Dougall Bell generator evaluates all its mixed combinations;
it does not show that those combinations span every conceivable
polygamma product. The unequal-twist scalar theorem concerns
identities of functions with two distinct branch centers;
it does not prove independence of their values on the integer
Fourier lattice. These qualifications are part of the results.

\subsection{Reproducibility and the role of computation}

The analytic proofs supply the infinite summations, domains,
continuations, and interchanges of limits. Four scripts implement
independent checks or exact finite algebra. Their JSON records
are included in \src{results/}; they contain the actual inputs,
working precision, and residuals. The scripts need Python,
\src{mpmath}, and \src{sympy}. The exact environment versions
and commands are recorded in the package's \src{README.txt}.

\begin{longtable}{@{}>{\raggedright\arraybackslash}p{0.28\textwidth}>{\raggedright\arraybackslash}p{0.63\textwidth}@{}}
\toprule
Artifact & Checks and interpretation\\
\midrule
\endhead
\src{verify_harmonic.py}
& Twenty-six numerical checks: twelve truncated Newton spectral
derivatives against Hurwitz derivatives, one complex-parameter
Mellin comparison, nine finite integer-shift specializations,
and four repeated antiderivatives. The Newton series uses
$240$ terms and $180$ working digits; its observed discrepancies
are about $10^{-18}$ to $10^{-15}$, set by truncation rather
than the working precision.\\[6pt]
\src{verify_dougall.py}
& Thirteen numerical comparisons of the trigamma square,
two-parameter Gamma identity, mixed spectral generator, and first
spectral derivative. Six exact Bell-polynomial checks and
twenty-one exact residue-grid factorizations use rational
symbolic arithmetic. The numerical residuals are below
$5.4\times10^{-46}$ in the supplied run.\\[6pt]
\src{verify_twists.py}
& Forty-nine exact partial-fraction checks through $p,q=7$,
one deliberately corrupted-sign rejection, and the exact
lowest-jet polarization identity. Eighteen numerical comparisons
cover ordinary Green convolutions, confluent twists, the bounded
logarithm difference, rational digamma conversion, branch
crossings, Fourier coefficients, and complex-order tail modes.\\[6pt]
\src{verify_collisions.py}
& Seventeen separated-collision diagnostics at $70$ working
digits, plus a two-split evaluation of the merged integral.
The independent interval quadratures visibly subtract each
simple pole. Their residuals test the predicted asymptotic
expansion; they are not exact-zero tests.\\
\bottomrule
\end{longtable}

The finite symbolic checks do not formalize the analytic proofs.
The floating-point comparisons are not interval enclosures.
In particular, an estimated asymptotic tail term is not a
certified error bound for the Gamma-ratio computations.
The harmonic finite differences can lose many digits through
cancellation; a larger working precision alone does not remove
the separate Newton truncation error. These are concrete
limitations of the supplied calculations, not qualifications
of the proved normal-convergence theorems.

\subsection{Further questions with concrete first targets}

The following are proposals for additional exact research. They
are not included in the proved theorem list.

\begin{question}[Which individual polygamma moments can be separated?]
At a fixed total order $w=r+s$, determine the exact span of the
mixed Bell polynomials $\mathcal B_{rs}$, allowing the elementary
linear polygamma rows and explicit parameter symmetries. The
square identity used the difference of two rows at $w=4$.
At $w=6$ and $w=8$, give exact rational elimination matrices and
identify which individual moments remain outside that span.
An absence from this specified span should be reported as such,
not as an impossibility of evaluation by other identities.
\end{question}

\begin{question}[Compact formulas at every integer residue intersection]
Use \eqref{dg:grid-first-mixed} to produce short explicit
tables at $(b,c)=(1,2)$ and $(2,2)$ for $N=2,3$.
Can central factorial polynomials organize the subtraction
coefficients without expanding large Bernoulli polynomials?
A useful certificate would retain the finite small-index
zeros and replay the mixed derivative using only exact
polynomial operations before evaluating polygamma endpoints.
\end{question}

\begin{question}[Finite reductions of regular harmonic spectral jets]
The Newton formula supplies an exact convergent representation;
classify parameter configurations at which it reduces to finitely
many ordinary Hurwitz-zeta order derivatives. Positive integral
$a$ and nonnegative integral $u$ give starting cases. For
positive depth and nonintegral rational $a$, first examine the
first two derivatives at $s=0,-1$. State the permitted function
class before asking whether a multiple Arakawa--Kaneko order
derivative is needed.
\end{question}

\begin{question}[The exact boundary of shift continuation]
The proved depth domain is $\Re a>-r$, and the Taylor series
about $a=1$ converges at least for $|a-1|<r+1$.
Determine the actual singular germ at $a=-r$ for generic complex
$s$ and compare it with the exceptional nonpositive integral
$s$, where the finite Stirling formula is polynomial in $a$.
This should distinguish a sufficient convergence domain from
a maximal analytic continuation and explain any exceptional
cancellation of the first shift singularity.
\end{question}

\begin{question}[Rational-shift character cancellations]
For rational $a$, transform the full mixed spectral germ
\eqref{dg:all-Stieltjes} into finite residue-class or
Dirichlet-character coordinates. The first target is $a=1/2$,
where $\gamma_m(a/2)=\gamma_m(1/4)$ appears. At each prescribed
order, compute an exact reduction matrix including every
conductor logarithm. Which parameter symmetries cancel a
whole character sector before numerical evaluation?
\end{question}

\begin{question}[Unequal twists approaching an integer resonance]
The finite-mode decomposition in \eqref{tw:tailseries} gives
the exact modes that change phase. Let one or both twists
approach an integer while the spectral orders are also
differentiated. After isolating the resonant Fourier modes,
determine a holomorphic completed germ and its endpoint contact
terms. The first target is a simultaneous pair approaching
zero at distinct fixed rates. The normalization should be
compared with the already known single-twist degeneration.
\end{question}

\begin{question}[From scalar functional reduction to sampled identities]
The finite scalar criterion proved here uses local monodromy.
Which additional hypotheses would let one infer a scalar
identity from equality on every integer Fourier mode in the
same finite power--logarithm class? An exact asymptotic uniqueness
theorem at infinity is one possible route. Conversely, seek
explicit sampled identities that do not extend as scalar
functional identities. This is distinct from any question of
arithmetic independence of individual special values.
\end{question}

\begin{question}[A complete geometric collision calculus]
Extend Theorem~\ref{cl:collision} to arbitrary Stieltjes indices
and argument-derivative orders, first for three unit-frequency
factors. The singular pieces are explicit derivatives of
$x^{-1}\log^m x$, so finite partial fractions and order
derivatives suggest an algorithm. It should output every
power--logarithm term and the finite rate-dependent term, with
an integrable uniform remainder. The desired result is an
identity for a specified normalization, not a sharper bound.
\end{question}

\begin{question}[Hierarchical collisions and composition]
Compare the paths $(0,\epsilon,\epsilon+\epsilon^2)$ and
$(0,\epsilon,2\epsilon)$ with iterated merging of a selected
pair followed by the third point. Can the finite corrections
be organized by nested subsets so that renormalized collision
composition is associative? The rate dependence already proved
shows why plain constant extraction is insufficient. Specify
the intermediate coordinate and subtraction at every stage.
\end{question}

\begin{question}[Permutation cancellations beyond total derivatives]
The completed generators distinguish global Tornheim data from
local endpoint data. Determine which nontrivial permutation
or character projections remove the global part for four
factors, after the elementary total-derivative identities have
been accounted for. A concrete target is total argument order
three with equal Stieltjes indices, retaining the complete
resonant-subset numerators. Produce either an exact cancellation
or a precise surviving multiple-sum kernel.
\end{question}

\begin{question}[Finite global kernels for four colliding grids]
Theorem~\ref{cl:allfactor} solves the local completion for any
number of factors. Construct a useful global spectral realization
for four frequencies on the rank-three lattice
$\sum_i p_i n_i=0$, with a finite sign-region and character
decomposition. The first target is $(p_1,p_2,p_3,p_4)=(1,1,2,3)$.
Local existence alone should not be presented as a minimal
finite reduction of the global multiple Tornheim kernels.
\end{question}

\begin{question}[The frozen Gaussian candidates]
Retain the exact source vectors for $S_6$ and the current $S_8$.
Seek an additional functional or iterated-integral relation
beyond the projection mechanisms that already prove $S_4$.
A successful proof should provide a finite exact row certificate
and an analytic justification for every generating row,
including its endpoint evaluation. An exclusion from a chosen
finite relation space would identify a missing mechanism,
not disprove the numerical identity.
\end{question}

\subsection{Recommended integration order}

The four mathematical sections can be integrated independently after
their cited antecedents. Preserve the stated domains, branches, and
normalizations, and review the two source edits separately.
The accompanying integration notes map labels and source targets
while preserving the existing conjecture status.

% END INLINED SOURCE: sections/06_audit_research.tex
\clearpage
% BEGIN INLINED SOURCE: references.tex
\begin{thebibliography}{99}
\raggedright
\addcontentsline{toc}{section}{References}

\bibitem{RepoMain}
V. Reshetnikov (repository maintainer),
\emph{ProveIt: Analysis/Polylogarithms/docs/manuscript}.
Snapshot \src{dcd95baeb7c0e914b138bddef6829c5b1d52b726},
10 October 2026.
\url{https://github.com/VladimirReshetnikov/ProveIt/tree/dcd95baeb7c0e914b138bddef6829c5b1d52b726/Analysis/Polylogarithms/docs/manuscript}.

\bibitem{RepoEndpoint}
ProveIt incoming research archive,
\src{ProveIt_Endpoint_Regularization_2026-10-10.zip},
especially \src{article/endpoint_regularization.tex}.
In \src{docs/incoming} at the snapshot in \cite{RepoMain}.

\bibitem{RepoDougall}
ProveIt incoming research archive,
\src{ProveIt_Dougall_Resonance_2026-10-10.zip},
centered Gamma-ratio resonance and spectral derivatives.
In \src{docs/incoming} at the snapshot in \cite{RepoMain}.

\bibitem{RepoTransport}
ProveIt incoming research archive,
\src{ProveIt_Resonance_Coordinate_Transport_2026-10-10.zip},
especially \src{sections/01_dougall.tex},
\src{sections/02_dougall_jets.tex}, and
\src{sections/08_research_questions.tex}, Question~3.
In \src{docs/incoming} at the snapshot in \cite{RepoMain}.

\bibitem{RepoExact}
ProveIt incoming research archive,
\src{ProveIt_Exact_Identities_2026-10-10.zip},
especially \src{sections/fully_shifted_height_one.tex},
\src{sections/collision.tex}, \src{sections/triple_dilation.tex},
and \src{sections/audit_research.tex}.
In \src{docs/incoming} at the snapshot in \cite{RepoMain}.

\bibitem{RepoCoincident}
ProveIt incoming research archive,
\src{ProveIt_Coincident_Stieltjes_2026-10-10.zip},
especially \src{article/article.tex}.
In \src{docs/incoming} at the snapshot in \cite{RepoMain}.

\bibitem{RepoTwisted}
ProveIt incoming research archive,
\src{ProveIt_Twisted_Stieltjes_Harmonic_Laurent_2026-10-10.zip},
especially \src{sections/03_twisted.tex},
\src{sections/04_harmonic_laurent.tex}, and
\src{sections/07_verification_questions.tex}.
In \src{docs/incoming} at the snapshot in \cite{RepoMain}.
The incoming directory is
\url{https://github.com/VladimirReshetnikov/ProveIt/tree/dcd95baeb7c0e914b138bddef6829c5b1d52b726/docs/incoming}.

\bibitem{DLMFDougall}
NIST Digital Library of Mathematical Functions,
\S16.4, especially equation 16.4.9, Rogers--Dougall summation.
\url{https://dlmf.nist.gov/16.4}.

\bibitem{DLMFGammaAsymptotic}
NIST Digital Library of Mathematical Functions,
\S5.11, especially equations 5.11.8 and 5.11.14,
Stirling--Bernoulli expansions and centered Gamma ratios.
\url{https://dlmf.nist.gov/5.11}.

\bibitem{DLMFPolygamma}
NIST Digital Library of Mathematical Functions,
\S5.15, polygamma functions.
\url{https://dlmf.nist.gov/5.15}.

\bibitem{DLMFHurwitz}
NIST Digital Library of Mathematical Functions,
\S25.11, Hurwitz zeta function.
\url{https://dlmf.nist.gov/25.11}.

\bibitem{DLMFPolylog}
NIST Digital Library of Mathematical Functions,
\S25.12, polylogarithms, especially equation 25.12.12.
\url{https://dlmf.nist.gov/25.12}.

\bibitem{DLMFLerch}
NIST Digital Library of Mathematical Functions,
\S25.14, Lerch transcendent.
\url{https://dlmf.nist.gov/25.14}.

\bibitem{DLMFHypergeom}
NIST Digital Library of Mathematical Functions,
\S15.8, transformations of the Gauss hypergeometric function.
\url{https://dlmf.nist.gov/15.8}.

\bibitem{AK1999}
T. Arakawa and M. Kaneko,
Multiple zeta values, poly-Bernoulli numbers, and related zeta functions,
\emph{Nagoya Mathematical Journal} \textbf{153} (1999), 189--209.
\href{https://doi.org/10.1017/S0027763000006954}{doi:10.1017/S0027763000006954}.
Author-hosted text:
\url{https://www2.math.kyushu-u.ac.jp/~mkaneko/papers/18MZV_pB_Nagoya.pdf}.

\bibitem{KT2018}
M. Kaneko and H. Tsumura,
Multi-poly-Bernoulli numbers and related zeta functions,
\emph{Nagoya Mathematical Journal} \textbf{232} (2018), 19--54.
Author-hosted text:
\url{https://www2.math.kyushu-u.ac.jp/~mkaneko/papers/56Kaneko_Tsumura.pdf}.

\bibitem{CC2010}
M.-A. Coppo and B. Candelpergher,
The Arakawa--Kaneko zeta function,
\emph{The Ramanujan Journal} \textbf{22} (2010), 153--162.
\href{https://doi.org/10.1007/s11139-009-9205-x}{doi:10.1007/s11139-009-9205-x}.
The exact Hurwitz normalization is on p.~154.
Author-hosted text:
\url{https://math.univ-cotedazur.fr/u/coppo/RamanujanJ(2010).pdf}.

\bibitem{EspinosaMoll}
O. Espinosa and V. H. Moll,
The evaluation of Tornheim double sums. Part~1,
\emph{Journal of Number Theory} \textbf{116} (2006), 200--229.
Preprint \src{arXiv:math/0505647}.
\url{https://arxiv.org/abs/math/0505647}.

\end{thebibliography}

% END INLINED SOURCE: references.tex
\end{document}
